% concatenated from preprint_main.tex
%%%%%%% ICML 2026 EXAMPLE LATEX SUBMISSION FILE %%%%%%%%%%%%%%%%%

\documentclass{article}

% Recommended, but optional, packages for figures and better typesetting:
\usepackage{microtype}
\usepackage{graphicx}
\usepackage{subfigure}
\usepackage{booktabs} % for professional tables
\usepackage{makecell} % for the appendix results table
% hyperref makes hyperlinks in the resulting PDF.
% If your build breaks (sometimes temporarily if a hyperlink spans a page)
% please comment out the following usepackage line and replace
% \usepackage{icml2026} with \usepackage[nohyperref]{icml2026} above.
\usepackage{hyperref}
\usepackage{graphicx}
\usepackage{rotating}

% some comment related packages 
\usepackage[textsize=tiny]{todonotes}
\usepackage{comment}




% Attempt to make hyperref and algorithmic work together better:
\newcommand{\theHalgorithm}{\arabic{algorithm}}

% Use the following line for the initial blind version submitted for review:
% \usepackage{Styles/icml2026}

% If accepted, instead use the following line for the camera-ready submission:
\usepackage[accepted]{Styles/icml2026}

% For theorems and such
\usepackage{amsmath}
\usepackage{amssymb}
\usepackage{mathtools}
\usepackage{amsthm}

% if you use cleveref..
\usepackage[capitalize,noabbrev]{cleveref}

%%%%%%%%%%%%%%%%%%%%%%%%%%%%%%%%
% THEOREMS
%%%%%%%%%%%%%%%%%%%%%%%%%%%%%%%%
\theoremstyle{plain}
\newtheorem{theorem}{Theorem}[section]
\newtheorem{proposition}[theorem]{Proposition}
\newtheorem{lemma}[theorem]{Lemma}
\newtheorem{corollary}[theorem]{Corollary}
\newtheorem{remark}[theorem]{Remark} % NOTE changed
\newtheorem{example}[theorem]{Example} % NOTE changed
\theoremstyle{definition}
\newtheorem{definition}[theorem]{Definition}
\newtheorem{assumption}[theorem]{Assumption}
\theoremstyle{remark}


% Todonotes is useful during development; simply uncomment the next line
%    and comment out the line below the next line to turn off comments
%\usepackage[disable,textsize=tiny]{todonotes}
\usepackage[textsize=tiny]{todonotes}



%%%%%%%%%%%%%%%%%%%%%%%%%%%%%%%%%%%%%%%%%%%%%%%%%%%%%%%%%%%%%%%%%
%%%%%%%%%%%%%%%%%%%%% START MY OWN IMPORTS %%%%%%%%%%%%%%%%%%%%%%
%%%%%%%%%%%%%%%%%%%%%%%%%%%%%%%%%%%%%%%%%%%%%%%%%%%%%%%%%%%%%%%%%

\usepackage{mathrsfs}
\usepackage{listings}
\usepackage{xcolor}

% Symbols
\newcommand{\bbR}{\mathbb{R}}
\newcommand{\bfR}{\mathbf{R}}
\newcommand{\scrR}{\mathscr{R}}
\newcommand{\calR}{\mathcal{R}}
\newcommand{\frakR}{\mathfrak{R}}
\newcommand{\bbQ}{\mathbb{Q}}
\newcommand{\bbP}{\mathbb{P}}

\newcommand{\bbS}{\mathbb{S}}
\newcommand{\bfS}{\mathbf{S}}
\newcommand{\scrS}{\mathscr{S}}
\newcommand{\calS}{\mathcal{S}}
\newcommand{\frakS}{\mathfrak{S}}
\newcommand{\bbX}{\mathbb{X}}
\newcommand{\bfX}{\mathbf{X}}
\newcommand{\scrX}{\mathscr{X}}
\newcommand{\calX}{\mathcal{X}}
\newcommand{\frakX}{\mathfrak{X}}
\newcommand{\bbY}{\mathbb{Y}}
\newcommand{\bfY}{\mathbf{Y}}
\newcommand{\bbZ}{\mathbb{Z}}
\newcommand{\bfZ}{\mathbf{Z}}
\newcommand{\scrF}{\mathscr{F}}
\newcommand{\calF}{\mathcal{F}}
\newcommand{\scrG}{\mathscr{G}}
\newcommand{\calG}{\mathcal{G}}
\newcommand{\scrA}{\mathscr{A}}
\newcommand{\calA}{\mathcal{A}}
\newcommand{\scrL}{\mathscr{L}}
\newcommand{\calL}{\mathcal{L}}
\newcommand{\bbO}{\mathbb{O}}
\newcommand{\bfO}{\mathbf{O}}
\newcommand{\scrO}{\mathscr{O}}
\newcommand{\calO}{\mathcal{O}}
\newcommand{\frakO}{\mathfrak{O}}
\newcommand{\scrD}{\mathscr{D}}
\newcommand{\calD}{\mathcal{D}}
\newcommand{\bbC}{\mathbb{C}}
\newcommand{\bfC}{\mathbf{C}}
\newcommand{\scrC}{\mathscr{C}}
\newcommand{\calC}{\mathcal{C}}
\newcommand{\frakC}{\mathfrak{C}}
\newcommand{\scrM}{\mathscr{M}}
\newcommand{\calM}{\mathcal{M}}
\newcommand{\frakM}{\mathfrak{M}}
\newcommand{\frakm}{\mathfrak{m}}
\newcommand{\bbT}{\mathbb{T}}
\newcommand{\bfT}{\mathbf{T}}
\newcommand{\scrT}{\mathscr{T}}
\newcommand{\calT}{\mathcal{T}}
\newcommand{\frakT}{\mathfrak{T}}
\newcommand{\bbK}{\mathbb{K}}
\newcommand{\bfK}{\mathbf{K}}
\newcommand{\scrK}{\mathscr{K}}
\newcommand{\calK}{\mathcal{K}}
\newcommand{\frakK}{\mathfrak{K}}
\newcommand{\bbB}{\mathbb{B}}
\newcommand{\bfB}{\mathbf{B}}
\newcommand{\scrB}{\mathscr{B}}
\newcommand{\calB}{\mathcal{B}}
\newcommand{\frakB}{\mathfrak{B}}
\newcommand{\bbH}{\mathbb{H}}
\newcommand{\bfH}{\mathbf{H}}
\newcommand{\scrH}{\mathscr{H}}
\newcommand{\calH}{\mathcal{H}}
\newcommand{\frakH}{\mathfrak{H}}
\newcommand{\bbW}{\mathbb{W}}
\newcommand{\bfW}{\mathbf{W}}
\newcommand{\scrW}{\mathscr{W}}
\newcommand{\calW}{\mathcal{W}}
\newcommand{\frakW}{\mathfrak{W}}

%% math operators
\DeclareMathOperator*{\argmax}{argmax}
\DeclareMathOperator*{\Ad}{Ad}
\DeclareMathOperator*{\argmin}{argmin}
\DeclareMathOperator*{\as}{a.s.}
\DeclareMathOperator*{\almosteverywhere}{a.s.}
\DeclareMathOperator{\sig}{\sigma-algebra}
\DeclareMathOperator*{\Prob}{Prob}
\DeclareMathOperator{\sgn}{sgn}
\DeclarePairedDelimiterX{\normop}[1]{\lVert}{\rVert_{\mathrm{op}}}{#1}
\DeclarePairedDelimiterX{\normopp}[1]{\lVert}{\rVert_{\mathrm{op}}^p}{#1}
\DeclarePairedDelimiterX{\normhs}[1]{\lVert}{\rVert_{\mathrm{HS}}}{#1}
\newcommand{\normtwo}[1]{\big\lVert#1\big\rVert_{2}}
\DeclarePairedDelimiter\ceil{\lceil}{\rceil}
\DeclarePairedDelimiter\floor{\lfloor}{\rfloor}

%% commands
\newcommand{\indicator}{\mathds{1}}
\newcommand{\E}{\mathbb{E}}
\newcommand{\var}{\mathrm{Var}}
\newcommand{\cov}{\mathrm{Cov}}
\newcommand{\tr}[1]{\mathrm{tr} \Big(#1\Big)}
\newcommand{\NC}[1]{\mathrm{NC}_2\big(#1\big)}
\newcommand{\wt}[1]{\mathrm{wt} (#1)}
\newcommand{\ubar}[1]{\underline{#1}}
\newcommand{\id}{\mathbbm{1}}
\newcommand{\notimplies}{\;\not\!\!\!\implies}
\newcommand\pvar[1]{\left\langle#1\right\rangle}
\newcommand\inner[2]{\langle#1,#2\rangle}
\def\delequal{\mathrel{\ensurestackMath{\stackon[1pt]{=}{\scriptscriptstyle A}}}}

\newcommand\smallO{
  \mathchoice
    {{\scriptstyle\mathcal{O}}}% \displaystyle
    {{\scriptstyle\mathcal{O}}}% \textstyle
    {{\scriptscriptstyle\mathcal{O}}}% \scriptstyle
    {\scalebox{.6}{$\scriptscriptstyle\mathcal{O}$}}%\scriptscriptstyle
  }
\newcommand{\bigO}{\calO}
%%%%%%%%%%%%%%%%%%%%%%%%%%%%%%%%%%%%%%%%%%%%%%%%%%%%%%%%%%%%%%%%%
%%%%%%%%%%%%%%%%%%%%%  END MY OWN IMPORTS  %%%%%%%%%%%%%%%%%%%%%%
%%%%%%%%%%%%%%%%%%%%%%%%%%%%%%%%%%%%%%%%%%%%%%%%%%%%%%%%%%%%%%%%%







% The \icmltitle you define below is probably too long as a header.
% Therefore, a short form for the running title is supplied here:
% \icmltitlerunning{Global Convergence Rates for Newton Boosting Trees}
\icmltitlerunning{Gradient Regularized Newton Boosting Trees with Global Convergence}
% TLDR: We formulate tree boosting as convex Restricted Newton Descent on Hilbert spaces, first analyzing vanilla Newton boosting, and then establishing a global O(1/k^2) convergence rate for a gradient-regularized boosting scheme.

\begin{document}

\twocolumn[
% \icmltitle{Global Convergence Rates for Newton Boosting Trees}
\icmltitle{Gradient Regularized Newton Boosting Trees with Global Convergence}

% It is OKAY to include author information, even for blind
% submissions: the style file will automatically remove it for you
% unless you've provided the [accepted] option to the icml2026
% package.

% List of affiliations: The first argument should be a (short)
% identifier you will use later to specify author affiliations
% Academic affiliations should list Department, University, City, Region, Country
% Industry affiliations should list Company, City, Region, Country

% You can specify symbols, otherwise they are numbered in order.
% Ideally, you should not use this facility. Affiliations will be numbered
% in order of appearance and this is the preferred way.

%\icmlsetsymbol{equal}{*}
\icmlsetsymbol{last}{$\dagger$}

\begin{icmlauthorlist}
\icmlauthor{Nikita Zozoulenko}{imperial}
\icmlauthor{Daniel Falkowski}{stockholm}
\icmlauthor{Thomas Cass}{imperial,last}
\icmlauthor{Lukas Gonon}{imperial,gallen,last}
\end{icmlauthorlist}

\icmlaffiliation{imperial}{Department of Mathematics, Imperial College London, UK}
\icmlaffiliation{gallen}{School of Computer Science, University of St. Gallen, Switzerland}
\icmlaffiliation{stockholm}{Department of Mathematics, Stockholm University, Sweden}

\icmlcorrespondingauthor{Nikita Zozoulenko}{n.zozoulenko23@imperial.ac.uk}
%\icmlcorrespondingauthor{XXX}{email@email.com}
% You may provide any keywords that you
% find helpful for describing your paper; these are used to populate
% the "keywords" metadata in the PDF but will not be shown in the document
\icmlkeywords{Machine Learning, ICML, gradient boosting, decision trees, convex optimization, Newton, xgboost, lightgbm, Newton boosting, Hilbert space}

\vskip 0.3in
]

% this must go after the closing bracket ] following \twocolumn[ ...

% This command actually creates the footnote in the first column
% listing the affiliations and the copyright notice.
% The command takes one argument, which is text to display at the start of the footnote.
% The \icmlEqualContribution command is standard text for equal contribution.
% Remove it (just {}) if you do not need this facility.

\printAffiliationsAndNotice{$\dagger$Equal last author}  % leave blank if no need to mention equal contribution
% \printAffiliationsAndNotice{\icmlEqualContribution} % otherwise use the standard text.
% \printAffiliationsAndNotice{$^\dagger$Equal last authors.}
%\printAffiliationsAndNotice{\icmlEqualContribution\hspace{-2pt}. $^\dagger$Equal last authors.}


\begin{abstract}
    Gradient Boosting Decision Trees (GBDTs) dominate tabular machine learning, with modern implementations like XGBoost, LightGBM, and CatBoost being based on Newton boosting: a second-order descent step in the space of decision trees. Despite its empirical success, the global convergence of Newton boosting is poorly understood compared to first-order boosting. In this paper, we introduce \textbf{Restricted Newton Descent}, which studies convex optimization with Newton's method on Hilbert spaces with inexact iterates, based on the concepts of cosine angle and weak gradient edge. Within this framework, we recover Newton boosting with GBDTs and classical finite-dimensional theory as special cases. 
    We first prove that vanilla Newton boosting achieves a linear rate of convergence for smooth, strongly convex losses that satisfy a Hessian-dominance condition. To handle general convex losses with Lipschitz Hessians, we extend a recent gradient regularized Newton scheme to the restricted weak learner setting. This scheme minimally modifies the classical algorithm by introducing an adaptive $\ell_2$-regularization term proportional to the square root of the gradient norm at each iteration. We establish a $\mathcal{O}(\frac{1}{k^2})$ rate for this scheme, thereby obtaining a globally convergent second-order GBDT algorithm with a rate matching that of first-order boosting with Nesterov momentum.
    In numerical experiments, we show that our scheme converges while vanilla Newton boosting may diverge.
\end{abstract}


\vspace{-20pt}
\section{Introduction}\label{sectionIntroduction}

Gradient Boosting Decision Trees (GBDTs) are powerful machine learning models, excelling in tabular settings such as regression, classification, and learning to rank over neural baselines \cite{
2021AreNeuralRankersStillOutperformedByGradientBoostedDecisionTrees,
2022WhyDoTreesOutperformDeepLearningTabularData,
2023WhenDoNeuralNetsOutperformBoostedTreesOnTabularData,
2025TabArenaBenchmarkTabularData}. Modern implementations such as XGBoost \cite{2016XGBoost}, LightGBM \cite{2017LightGBM}, and CatBoost \cite{2018CatBoost} are based on the Newton boosting algorithm, which updates a decision tree ensemble by performing a Newton descent step in the space of decision trees, derived via a second order Taylor expansion with respect to the training data \cite{2016XGBoost}. However, despite the widespread use of GBDTs, the theoretical mechanisms of Newton boosting remain poorly understood, with most prior work focused on first-order gradient boosting \cite{1999BoostingAlgorithmsAsGradientDescent,
2001GreedyFunctionApproximationAGradientBoostingMachine,
2011GeneralizedBoostingAlgorithmsForConvexOptimization,
2017ANewPerspectiveOnBoostingInLinearRegressionViaSubgradientOptimization,
2019RegularizedGradientBoosting,
2020RandomizedGradientBoostingMachine,
2020AcceleratingGradientBoostingMachines}.


While Newton's method in $\mathbb{R}^N$ offers fast quadratic local convergence, it lacks global convergence guarantees for general strictly convex losses, even when combined with line search \cite{2016SimpleExamplesForTheFailureOfNewtonsMethodWithLineSearchForStrictlyConvexMinimization}. Recently, the convex optimization literature has seen a resurgence of interest in regularized Newton schemes that achieve global $\mathcal{O}(\frac{1}{k^2})$ convergence rates. This was first established for the Cubic Regularized Newton (CRN) method \cite{2006CubicRegularizationOfNewtonMethodAndItsGlobalPerformance, 2008AcceleratingTheCubicRegularizationOfNewtonsMethodOnConvexProblems}; however, CRN requires solving an expensive subproblem at each iteration, making it unfit for GBDTs. More recently, the Gradient Regularized Newton (GRN) method was proposed, achieving the same $\mathcal{O}(\frac{1}{k^2})$ rate as CRN \cite{2023RegularizedNewtonMethodWithGlobalConvergence, 2024GradientRegularizationOfNewtonMethodWithBregmanDistances}. GRN uses an adaptive $\ell_2$-regularization term proportional to the square root of the gradient norm at each iteration. Crucially, this method avoids expensive subproblems and incurs negligible computational overhead compared to the standard Newton update. This discovery has sparked new research into the global convergence of second-order numerical schemes \cite{2022ADampedNewtonMethodAchievesGlobalMinus2AndLocalQuadraticConvergenceRate, 2024SuperUniversalRegularizedNewtonMethod, 2025MinimizingQuasiSelfConcordantFunctionsByGradientRegularizationOfNewtonMethod, 2025SketchAndProjectMeetsNewtonMethodGlobalConvergenceWithLowRankUpdates}. A natural question arises: can these convergence guarantees be extended to GBDTs and optimization with weak learners, providing convergence rates for second-order boosting methods used in practice?
 		







% Newton's method in a functional space of trees, most prior theoretical work has focused on variants of the original formulation of friedman et al which performs a gradient descent step, rather than a Newton step for the tree ensemble update (TODO CITE). 


% gradient boosting models have traditionally been viewed as coordinate or gradient descent in $\mathbb{R}^n$ on the training data, where $n$ is dataset size. Later, this has been identified with the hilbert space of square integrable $L^2(P_n)$ where $P_n$ is the empirical distribution of the data, see also (CITE TODO for connections of boosting to neural networks).

% recently works have taken to study gradient boosting via (randomized) coordinate descent. This can be done by noting that tree based models are limited by the quantiles of the training data, hence can be enumerated by a finite family of vectors, leading to the coordinate descent point of view. In this paper, we argue that it may be more natural and elegant to simply view boosting as convex optimization in hilbert spaces, where optimization is done using weak iterates following a weak learner edge or cosine angle assumption. This way, both tree boosting and the classical finite dimensional convex optimization results can be optained as special cases.

\subsection{Contributions}
In this paper, we formulate Newton boosting as a convex optimization problem on a Hilbert space $\mathcal{H}$, where optimization is performed on a restricted set of permissible weak learners $\mathcal{F} \subset \mathcal{H}$. Our contributions are:

\begin{itemize}
    \item We introduce \textbf{Restricted Newton Descent} using weak iterates based on the concepts of cosine angle and weak gradient edge, obtaining both tree boosting and finite-dimensional convex optimization as special cases.

    \item For Hessian-dominated smooth strongly convex losses, we prove that vanilla restricted Newton's method achieves a linear rate of convergence. This includes $\ell_2$-regularized binary and categorical cross entropy.

    \item For general convex losses with $2M$-Lipschitz Hessians, we generalize the gradient regularized Newton scheme of \citet{2023RegularizedNewtonMethodWithGlobalConvergence} to the Hilbert space setting with weak learners. We prove a global $\mathcal{O}(\frac{1}{k^2})$ convergence rate for all learning rates $\eta \leq 1$, and a local linear rate if the loss additionally is strongly convex.
    
    \item We present numerical results showing that the gradient regularized Newton scheme converges while vanilla Newton boosting may diverge.    
\end{itemize}


\subsection{Related Literature}


\textbf{GBDTs:} Gradient boosting models have traditionally been viewed as gradient descent in $\mathbb{R}^N$ on the training data, where $N$ is the dataset size \cite{1999BoostingAlgorithmsAsGradientDescent, 2000LogitBoostAdditiveLogisticRegressionAStatisticalViewOfBoosting}. Later works identified this process with gradient descent in the Hilbert space of square-integrable functions $L^2(\widehat{\nu}_N)$, where $\widehat{\nu}_N$ is the empirical data distribution \cite{2011GeneralizedBoostingAlgorithmsForConvexOptimization}; see also connections to gradient representation boosting in deep learning \cite{2018FunctionalGradientBoostingResNet, 2020GeneralizedBoosting, 2025RandomFeatureRepresentationBoosting}. Recent theoretical work has shifted from the gradient perspective to studying boosting via coordinate descent in high-dimensional spaces, achieved by enumerating all possible decision trees \cite{2019RegularizedGradientBoosting, 2020RandomizedGradientBoostingMachine, 2020SnapBoostAHeterogeneousBoostingMachine}.



\textbf{Newton Boosting:} Newton-like updates were first introduced by LogitBoost \cite{2000LogitBoostAdditiveLogisticRegressionAStatisticalViewOfBoosting} for logistic loss, and later applied to GBDTs in a general context by XGBoost \cite{2016XGBoost}. Empirical differences between first-order and Newton boosting were investigated in \cite{2021GradientAndNewtonBoostingForClassificationAndRegression}. Theoretical analysis of Newton boosting remains limited compared to first-order methods; Notable work includes the LogitBoost analysis by \citet{2014AConvergenceRateAnalysisForLogitBoostMARTAndTheirVariant} using probability clamping, and later the SnapBoost work using coordinate descent \cite{2020SnapBoostAHeterogeneousBoostingMachine}.

\textbf{Convergence Rates in Boosting:} For second order boosting, both \citet{2014AConvergenceRateAnalysisForLogitBoostMARTAndTheirVariant} and \citet{2020SnapBoostAHeterogeneousBoostingMachine} proved a linear rate in the strongly convex setting. In contrast, our work also covers general convex losses with Lipschitz Hessians. As for first order boosting, the earliest explicit convergence rate was proved by \citet{2006SomeTheoryForGeneralizedBoostingAlgorithms} with a slow $\mathcal{O}((\log k)^{-\frac{1}{2}})$ rate. For smooth convex losses, \citet{2014AConvergenceRateAnalysisForLogitBoostMARTAndTheirVariant} proved a classical $\mathcal{O}(\frac{1}{k})$ rate.  \citet{2020AcceleratingGradientBoostingMachines} incorporated Nesterov momentum to achieve an accelerated $\mathcal{O}(\frac{1}{k^2})$ rate, similar to our work. Linear rates for strongly convex losses were proven by \citet{2011GeneralizedBoostingAlgorithmsForConvexOptimization} using gradient descent, and by \citet{2020RandomizedGradientBoostingMachine} via randomized coordinate descent.  \citet{2012APrimalDualConvergenceAnalysisOfBoosting} studied the primal-dual structure of boosting, obtaining a linear rate with non-standard constants.



%TODO ADD EXTRA SENTENCE: other recent theretical work not directly related to global convergence of GBDTs include TOOD XXXXX...








\section{Newton Boosting as Hilbert Space Convex Optimization}
In this section, we recall Newton's method in the Hilbert space setting, and introduce the \textbf{Restricted Newton Descent} framework where Newton boosting with GBDTs is recovered as a special case.

\subsection{Hilbert Space Newton's Method}

Let $\mathcal{H}$ be a Hilbert space. In this paper we consider twice Fréchet differentiable convex loss functions $L : \mathcal{H} \to \mathbb{R}$, defined below.

\begin{definition}
    A function $L : \mathcal{H} \to \bbR$ is said to be \textbf{Fréchet differentiable} if there exists a mapping $\nabla L : \mathcal{H} \to \mathcal{H}$ such that for all $h,f \in \mathcal{H}$
    \begin{align*}
        L(h + f) = L(h) + \langle f, \nabla L (h)\rangle + o(\|f\|).
    \end{align*}
    The element $\nabla L(h)$ is called the \textbf{gradient} at $h$. A function $L : \mathcal{H} \to \bbR$ is said to be \textbf{twice Fréchet differentiable} if $L$ is Fréchet differentiable, and if for all $h\in\mathcal{H}$ there exists a bounded, self-adjoint linear operator $\nabla^2 L(h) : \mathcal{H} \to \mathcal{H}$ such that
    \begin{align*}
    L(h + f) = L(h) + \left\langle f, \nabla L(h) + \frac{1}{2}  \nabla^2 L(h)[f]
   \right\rangle + o(\|f\|^2),
    \end{align*}
    for all $f \in \mathcal{H}$. The map $\nabla^2 L(h)$ is called the \textbf{Fréchet Hessian} (or second derivative) of $L$ at $h$.
\end{definition}

We aim to minimize $L$ with Hilbert space Newton's method with a fixed learning rate $\eta>0$. At each iteration $k$, we update an iterate $F_k = F_{k-1} + \eta f_k = \sum_{t=1}^k \eta f_t$ by minimizing a second-order Taylor expansion of the loss:
\begin{align*}
    L(F_k + f) &\approx L(F_k) + \left\langle f, \nabla L(F_k) + \frac{1}{2}\nabla^2 L(F_k)[f]  \right\rangle \\
    &=  L(F_k) + \left\langle f,g_k + \frac{1}{2}H_k[f]  \right\rangle \\
    &=: L(F_k) + Q_k(f)
\end{align*}
where $g_k = \nabla L(F_k)$ is the gradient, $H_k = \nabla^2 L(F_k)$ is the Hessian operator, and $Q_k(f)$ is the local second-order approximation of $L$ at the point $F_k$, that is after $k$ iter\-ations of the optimization scheme. Differentiating $Q_k$ yields
\begin{align*}
    \nabla Q_k(f) = g_k + H_k[f].
\end{align*}
Setting the derivative of $Q_k$ to 0 and solving for $f$ gives the Newton update, assuming that $H_k$ is invertible at the point $g_k$:
\begin{align}\label{eqExactNewtonUpdate}
    f_{k+1}  := - H_k^{-1}[g_k] = \argmin_{f \in \mathcal{H}} Q_k(f).
\end{align}




\subsection{Restricted Newton Descent}
In the context of gradient boosting, the exact update $f_{k+1}$ in \eqref{eqExactNewtonUpdate} is generally not attainable since weak learners, as opposed to exact Newton directions, are used to minimize the loss. Consequently, the ensemble $F_k$ instead takes the form $F_k := \sum_{t=1}^k \eta f_t^w$ for some notion of a weak iterate $f^w_t$. To formalize this, we introduce the \textbf{Restricted Newton Descent} framework (not to be confused with constrained optimization). This differs from standard Newton's method in two ways. Firstly, the gradient $g_k$ and Hessian $H_k$ are computed w.r.t. the sum of weak, rather than exact, iterates $F_k$. Secondly, we do not minimize $Q_k(f)$ over the entire space $f \in \mathcal{H}$, but rather restrict the weak update $f_{k+1}^w$ to belong to some subset $f_{k+1}^w \in \mathcal{F} \subset \mathcal{H}$ of admissible weak learners. More specifically, given a family $\mathcal{F}$, we define the weak iterates $f_{k+1}$ as any element
\begin{align}\label{eqWeakNewtonUpdate}
    f_{k+1}^w  \in  \argmin_{f \in \mathcal{F}} Q_k(f).
\end{align}
Note that $f_{k+1}^w$ is not necessarily unique (although it will be for decision trees). We consider general families $\mathcal{F}$ satisfying the following scalability condition, similar to \citet{2020AcceleratingGradientBoostingMachines}:
\begin{assumption}
    The family $\mathcal{F}$ is closed under scalar multiplication. That is, if $f \in \mathcal{F}$ then $a f \in \mathcal{F}$ for all $a\in\mathbb{R}$.
\end{assumption}
Decision trees, commonly used in boosting, satisfy this assumption. If $\mathcal{F} = \mathcal{H}$, then the weak update $f^w_{k+1}$ coincides with the exact update $f_{k+1}$, and standard Newton's method is recovered. The full algorithm is detailed in \cref{alg:boostingVanillaNewton}, with the exact step in parentheses.

\begin{algorithm}[H]
\caption{Vanilla Restricted Newton}
\label{alg:boostingVanillaNewton}
\begin{algorithmic}[1]
    \STATE \textbf{Input:} Loss $L:\mathcal{H}\to\bbR$. Learning rate $\eta$. Weak learner family $\mathcal{F}$. Initial value $F_0$.
    \FOR{$k = 0, 1, 2, \dots$}
        \STATE $g_k, \: H_k \leftarrow \nabla L(F_k), \: \nabla^2L(F_k)$
        \STATE $\big($Exact target $f_{k+1} \leftarrow - H_k^{-1}g_k \big)$
        \STATE Weak learner $f^{w}_{k+1} \leftarrow \argmin_{f \in \mathcal{F}} Q_k(f)$
        \STATE $F_{k+1} \leftarrow F_k + \eta f^{w}_{k+1}$
    \ENDFOR
\end{algorithmic}
\end{algorithm}

Theoretical analysis of any optimization scheme now depends on the error between the chosen weak iterate $f^w_{k+1}$ and the exact iterate $f_{k+1}$, as this difference introduces an error component orthogonal to the optimal descent direction. As we show later, different schemes rely on different error metrics. For vanilla Newton's method, the relevant metric is the cosine angle between $f_{k+1}$ and $f^w_{k+1}$ in the induced Hessian norm (see \cref{secVanillaNewtonBoosting}), whereas for the globally convergent gradient regularized scheme, the metric takes the form of a weak learner edge condition between the true gradient and the implied weak gradient (see \cref{secGradientRegularizedNewton}).

\subsection{GBDTs as Restricted Convex Optimization}\label{subsec:GBDTsAsRestrictedConvexOptimization}

In this section, we provide an overview of gradient boosting from a functional $\mathcal{L}^2(\nu)$ perspective \cite{2011GeneralizedBoostingAlgorithmsForConvexOptimization, 2020GeneralizedBoosting}, in the context of Restricted Newton Descent.

Let $(X, Y) \sim \nu$ be a random sample with its corresponding probability measure $\nu$. Denote by $X \sim \nu_X$ the features in $\bbR^q$, with targets $Y \sim \nu_Y$ in $\bbR^K$, and $\nu_{Y|X=x}$ the conditional law. Let $\widehat{\nu}_N = \frac{1}{N}\sum_{i=1}^N \delta_{(x_i, y_i)}$ be the empirical measure of $\nu$ with respect to a training dataset $\{x_i, y_i\}_{i=1}^N$ of size $N$. We work in the Hilbert space $\mathcal{L}^2(\nu_X)$ of square integrable functions defined on the feature space, with inner product $\langle f, g\rangle = \E_{\nu_X}[\langle f(X),g(X) \rangle_{\bbR^K}]$. When $\nu = \widehat{\nu}_N$, this space can be identified with $\mathbb{R}^{N\times K}$ using the inner product $\langle f, h\rangle = \frac{1}{N}\sum_{i=1}^N \langle f(x_i),h(x_i) \rangle_{\bbR^K}$. Let $l$ be a twice differentiable convex loss function.

To derive the Newton boosting algorithm used by e.g. XGBoost and LightGBM as a special case of Restricted Newton Descent, we simply let $\mathcal{F}$ be the set of weak decision trees up to a fixed depth, let the base space be $\mathcal{H} = \mathcal{L}^2(\widehat{\nu}_{X,N})$, and the loss function be $L(F) = \frac{1}{N} \sum_{i=1}^N l(F(x_i), y_i) = \E_{\nu_N}[l(F(X),Y)]$. In the space $\mathcal{L}^2(\nu_X)$, the Fréchet gradient and Hessian take the form
\begin{align}
    \nabla L (F)(x)  =  \E_{\nu_{Y|X=x}}\big[ \partial_{1} l(F(x), Y) \big],
\end{align}
\begin{align}
    \nabla^2 L (F)[f](x)  =  \E_{\nu_{Y|X=x}}\big[ \partial_{11} l(F(x), Y) f(x) \big],
\end{align}
and are easily computed for empirical measures $\widehat{\nu}_N$ as partial derivatives and matrices with respect to $F(x_i)$. Write $g_{(i)} = \partial_1 l(F_k(x_i), y_i) \in \bbR^K$ and $h_{(i)} = \partial_{11} l(F_t(x_i), y_i) \in \bbR^{K\times K}$. Following the XGBoost derivation \cite{2016XGBoost}, the objective becomes to minimize
\begin{align*}
        L(F_k+f) 
        &= \frac{1}{N}\sum_{i=1}^N l(F_k(x_i)+f(x_i), y_i) \\
        &\approx  L(F_k) + \frac{1}{N}\sum_{i=1}^N \left\langle f(x_i), g_{(i)} + \frac{1}{2} h_{(i)} f(x_i) \right\rangle_{\mathbb{R}^K} \\
        &= L(F_k) + Q_k(f)
\end{align*}
which we see coincides with the Hilbert space derivation using $Q_k$. Furthermore, in boosting we minimize over the set of weak learners $f \in \mathcal{F}$ rather than over all $L^2(\widehat{\nu}_N)$. Letting $f$ be a decision tree of the form $f(x) = \sum_j w_j I_j(x)$ where $I_j(x) \in \{0, 1\}$ is a disjoint partition of the feature space, one finds that the optimal weights $w_j$ minimizing $Q_k(f)$ are
\begin{align*}
    w_j &=  -\left( \sum_{i \in I_j} h_{(i)} \right)^{-1} \sum_{i \in I_j} g_{(i)},
\end{align*}
coinciding with the classical XGBoost formula. See \cref{appendixXGBFormulaDerivation} for more details and full derivation.



\section{Restricted Newton for Hessian-dominated Losses}\label{secVanillaNewtonBoosting}

Before introducing the globally convergent scheme in \cref{secGradientRegularizedNewton}, we first analyze the vanilla Newton update via Restricted Newton Descent. We prove for smooth strongly convex losses satisfying a Hessian-dominance condition, that vanilla Newton's method with weak learners converges globally with a linear rate. This analysis broadly parallels that of the LogitBoost analysis \cite{2014AConvergenceRateAnalysisForLogitBoostMARTAndTheirVariant}, but we analyze it within our Restricted Newton Descent formulation in a more general setting beyond logistic loss. To ensure the Newton update is well-defined on $\mathcal{H}$, we assume the following:

\begin{assumption}\label{assumptionHessianNewtonEquationUniqueSolution}
    We assume that the equation $g_k = -H_k f$ admits a unique solution $f=f_{k+1} = - H_k^{-1} g_k$.
\end{assumption}

\cref{assumptionHessianNewtonEquationUniqueSolution} is satisfied, for example, when $L$ is strongly convex, or in many practical settings when $\mathcal{H}$ is a finite dimensional space with a strictly convex loss $L$. Note that this assumption is not required for the gradient regularized scheme in \cref{secGradientRegularizedNewton}.



\subsection{Cosine Angle}

Recall that the iterates of Restricted Newton Descent are derived by minimizing $Q_k(f) = \left\langle f, g_k + \frac{1}{2}H_k f \right\rangle$ restricted to $f\in\mathcal{F}$ at each iteration $k$, yielding $F_k = \sum_{t=1}^k \eta f^w_t$. In order to analyze the restricted Newton scheme, we need a way to measure the error between the exact direction $f_{k+1}$ and the weak learner $f^w_{k+1}$. It turns out that the critical quantity determining the performance of this scheme, is the cosine angle between $f_{k+1}$ and $f^w_{k+1}$ measured in the Newton geometry induced by the Hessian. For this, we use the following norm definition.




\begin{definition}\label{defCosineAngle}
    Let $A : \mathcal{H} \to \mathcal{H}$ be a self-adjoint positive operator. We define the inner product induced by $A$ as
    \begin{equation}
        \langle h, f\rangle_{A} := \langle h, A f\rangle.
    \end{equation}
    When $A = H_k$, we call this the \textbf{Hessian-induced inner product}, with induced norm $\|\cdot\|_{H_k}$.
\end{definition}

The following lemma motivates the use of the Hessian-induced norm and naturally leads to the definition of the cosine angle between weak and exact iterates.

\begin{lemma}\label{lemmaVanillaNewtonClosestElement}
    Minimizing $Q_k(f)$ is equivalent to finding a closest element $f\in\mathcal{F}$ to the exact Newton direction $f_{k+1}$ in the Hessian-induced norm $\|\cdot\|_{H_k}$.
\end{lemma}
\begin{proof}
    Using the Hessian-induced norm and the exact Newton update $f_{k+1}$, we can write
    \begin{align*}
        \|f-f_{k+1}\|^2_{H_k}
        &= \|f\|_{H_k}^2 - 2\langle f, f_{k+1}\rangle_{H_k} + \| f_{k+1} \|_{H_k}^2 \\
        &= \langle f, H_k f\rangle + 2\langle f, g_k \rangle + \|f_{k+1}\|^2_{H_k} \\
        &= 2Q_k(f) + \|f_{k+1}\|^2_{H_k}. \qedhere
    \end{align*}
\end{proof}


\begin{definition}
    Consider the iterates produced by vanilla Newton's method. At each iteration $k$, we define the Hessian-induced cosine angle $\Theta_k \in [0, 1]$ between the exact update $f_{k+1}$ and the weak learner $f_{k+1}^w$ as
    \begin{equation}\label{eqHessianCosineAngle}
        \Theta_k := \frac{\langle f_{k+1}, f_{k+1}^w \rangle_{H_k}}{\|f_{k+1}\|_{H_k} \|f_{k+1}^w\|_{H_k}}, \qquad \Theta := \inf_k \Theta_k.
    \end{equation}
\end{definition}

In contrast to \citet{2014AConvergenceRateAnalysisForLogitBoostMARTAndTheirVariant}, which uses a weak learnability assumption specific to binary classification, more recent prior work uses definitions similar in spirit to ours, but defines them by taking suprema and infima over much larger spaces of permissible weak learners and targets \cite{2020RandomizedGradientBoostingMachine, 2020AcceleratingGradientBoostingMachines,2020SnapBoostAHeterogeneousBoostingMachine}; see also \cite{2011GeneralizedBoostingAlgorithmsForConvexOptimization}. Our definition is simpler and more direct, capturing the specific two vectors $f^w_{k+1}$ and $f_{k+1}$ determining the cosine angle. Using this definition, we immediately obtain the following properties connecting the weak and exact Newton iterates with the gradient $g_k$. 

\begin{lemma}\label{lemmaCosineAngleProperties}
    If $\mathcal{F}$ is closed under scalar multiplication, then
    \begin{enumerate}
        \item $\|f_{k+1} - f_{k+1}^w\|^2_{H_k} = (1-\Theta_k^2) \|f_{k+1}\|^2_{H_k}$,
        \item $\|f_{k+1}^w\|_{H_k} = \Theta_k \|f_{k+1}\|_{H_k}$,
        \item $ \|f^w_{k+1}\|_{H_k}^2 = -\langle g_k, f_{k+1}^w\rangle$.
    \end{enumerate}
\end{lemma}
\begin{proof}
    Since $\mathcal{F}$ is closed under scalar multiplication, the objective $Q_k(\alpha f_{k+1}^w)$ is minimized at $\alpha=1$. Differentiating the equivalent objective (see \cref{lemmaVanillaNewtonClosestElement}) $\phi(\alpha) := \frac{1}{2}\|\alpha f_{k+1}^w - f_{k+1}\|_{H_k}^2$ with respect to $\alpha$ gives the derivative $\phi'(\alpha) = \langle f_{k+1}^w, \alpha f_{k+1}^w - f_{k+1} \rangle_{H_k}$. Setting this to zero at $\alpha=1$ yields
    $
        \langle f_{k+1}, f_{k+1}^w \rangle_{H_k} = \|f_{k+1}^w\|_{H_k}^2.
    $
    Substituting this into the definition of $\Theta_k$, we obtain
    \begin{equation}
        \Theta_k = \frac{\|f_{k+1}^w\|_{H_k}^2}{\|f_{k+1}\|_{H_k} \|f_{k+1}^w\|_{H_k}} = \frac{\|f_{k+1}^w\|_{H_k}}{\|f_{k+1}\|_{H_k}},
    \end{equation}
    which proves the second identity. The first identity then follows immediately from expanding the norm using inner products. The third identity follows from the fact that $-\langle g_k, f_{k+1}^w\rangle = \langle H_k f_{k+1}, f_{k+1}^w\rangle = \langle f_{k+1}, f_{k+1}^w\rangle_{H_k}$.
\end{proof}


\subsection{Decrease in Loss for Vanilla Restricted Newton}

We now show that the restricted Newton scheme guarantees a strict decrease in loss at each iteration, forming the basis of our convergence analysis. We rely on the following standard result from optimization theory, see \cref{app:SmoothnessLemmas} for more details.

\begin{lemma}\label{lemmaStandardSmoothnessUpperBound}
    If the loss $L$ is $S$-smooth, then for all $f\in\mathcal{H}$
    \begin{align}
        L(F_k+f) \leq L(F_k) + \langle g_k, f\rangle + \frac{S}{2} \|f\|^2.
    \end{align}
\end{lemma}

\begin{definition}\label{defHessianDominatedLoss}
    We say that the loss $L$ is \textbf{Hessian-dominated} if there exists a constant $c>0$ such that the exact iterates $f_{k+1}$ satisfy $\|f_{k+1}\|^2_{H_k} \geq c L(F_k)$.
\end{definition}

\cref{defHessianDominatedLoss} generalizes the condition derived in \citep[Theorem 4]{2014AConvergenceRateAnalysisForLogitBoostMARTAndTheirVariant} for LogitBoost, to general losses using the language of Hessian norms and weak iterates. Examples of Hessian-dominated losses include binary and categorical cross entropy, both with constant $c=1$ (see \cref{app:HessianDominatedLoss} for more details).

We are now ready to state and prove our main convergence result for Restricted Newton Descent.

\begin{theorem}[\textbf{Global Rate I}]\label{lemmaStandardSmoothnessResults}
    If $L$ is $S$-smooth, $\mu$-strongly convex, and $\mathcal{F}$ closed under scalar multiplication, then for any learning rate $\eta \in (0, \frac{2\mu}{S})$ we have
    \begin{align}
        L(F_k) - L(F_{k+1}) 
        &\geq \eta\left( 1 -\frac{\eta S}{2\mu}\right) \Theta^2\|f_{k+1}\|_{H_k}^2.
    \end{align}
    Consequently, if $L$ satisfies \cref{defHessianDominatedLoss}, then vanilla restricted Newton's method converges linearly with rate $(1-\rho)$, where $\rho = c\Theta^2 \eta(1 - \frac{\eta S}{2\mu})$:
    \begin{align}
        L(F_{k+1}) \leq (1-\rho) L(F_k).
    \end{align}
\end{theorem}
\begin{proof}
    Starting with the standard smoothness upper bound $L(F+f) \leq L(F) + \langle g, f \rangle + \frac{S}{2}\|f\|^2$ from \cref{lemmaStandardSmoothnessUpperBound}, we substitute the weak update $f = \eta f^w_{k+1}$. Using strong convexity to write $\|u\|^2 \leq \frac{1}{\mu}\|u\|_{H_k}^2$, and the cosine angle properties from \cref{lemmaCosineAngleProperties}, we obtain
    \begin{align}
        L(F_k) -L(F_{k+1})
        &\geq - \eta \langle g_k, f_{k+1}^w \rangle - \frac{S \eta^2}{2} \|f_{k+1}^w\|^2 \\
        &\geq \eta \|f_{k+1}^w\|_{H_k}^2 - \frac{S \eta^2}{2\mu} \|f_{k+1}^w\|_{H_k}^2 \\
        &= \left(\eta -\frac{S \eta^2}{2\mu} \right)\Theta_k^2\|f_{k+1}\|_{H_k}^2.
    \end{align}
    Applying \cref{defHessianDominatedLoss}, we get $L(F_{k+1}) \leq L(F_k) - \rho L(F_k)$, which implies linear convergence.
\end{proof}
\begin{remark}
    Our bound compares favourably to SnapBoost \cite{2020SnapBoostAHeterogeneousBoostingMachine}, which proved a linear rate of $1-\rho_{snap} = 1 - \Theta^2_{snap}\frac{\mu^2}{S^2}$ for a fixed learning rate $\eta = \frac{\mu}{S}$. While our definitions of $\Theta$ are comparable, our result yields a better dependence on the condition number $\frac{S}{\mu}$. Furthermore, we obtain essentially the same rate as \citet{2014AConvergenceRateAnalysisForLogitBoostMARTAndTheirVariant}, where they instead use a change of Hessian lemma, and their probability clamping is replaced by strong convexity in our lemma.
\end{remark}

While cross entropy loss itself is not strongly convex, this issue is typically addressed in two ways. First, for practical applications $\ell_2$-regularization is almost always applied to the loss function, yielding a strongly convex loss. Second, probability values can be clipped to $[\epsilon, 1-\epsilon]$ for some $\epsilon>0$ as in \cite{2014AConvergenceRateAnalysisForLogitBoostMARTAndTheirVariant}, which implicitly implies strong convexity. An alternative approach to \cref{lemmaStandardSmoothnessResults} could be to clip the output tree values, in combination with a \textit{change of Hessian Lemma}, similar to \citep[Lemma 9]{2014AConvergenceRateAnalysisForLogitBoostMARTAndTheirVariant} to obtain a linear rate. However, we believe our approach offers a more concise and simpler theoretical treatment.


\section{Gradient Regularized Restricted Newton}\label{secGradientRegularizedNewton}
Despite the fast quadratic local convergence of Newton's method in finite-dimensional convex optimization, the method does not have global convergence guarantees for general smooth convex functions. A classical example of this is $L(x) = \sqrt{1+x^2}$ defined for $x\in \mathbb{R}$, see \cref{app:newton_divergence}. However, cubic regularized Newton \cite{2006CubicRegularizationOfNewtonMethodAndItsGlobalPerformance} and the more recent gradient regularized Newton \cite{2023RegularizedNewtonMethodWithGlobalConvergence,2024GradientRegularizationOfNewtonMethodWithBregmanDistances} achieve a global $O(\frac{1}{k^2})$ convergence rate. In this section, we extend the latter scheme to Newton boosting with weak learners within our framework of Restricted Newton Descent.

Instead of minimizing the local quadratic approximation $Q_k(f)$ at each iteration $k$, gradient regularized Newton's method minimizes the following regularized objective:
\begin{align*}
    Q_k^{reg}(f) 
    &= \left\langle f,g_k + \frac{1}{2}H_k[f]  \right\rangle + \frac{\lambda_k}{2}\|f\|^2 \\
    &= Q_k(f) + \frac{\lambda_k}{2}\|f\|^2,
\end{align*}
where $\lambda_k \in (0, \infty)$ is a scalar regularization parameter. For losses with $2M$-Lipschitz continuous Hessians, we will allow values $\lambda_k \geq \sqrt{M\|g_k\|}$. Minimizing $Q_k^{reg}(f)$ yields the exact regularized update
\begin{align}\label{eqRegBoostingRoundExact}
    f_{k+1}  = - (H_k + \lambda_k I)^{-1}[g_k].
\end{align}
In other words, gradient regularized Newton is simply an $\ell_2$-regularized Newton step where the $\ell_2$-regularization is adaptive at each iteration, being proportional to the square root of the gradient norm $\sqrt{\|g_k\|}$. Similar to \cref{secVanillaNewtonBoosting}, we consider Restricted Newton Descent using weak learners rather than the exact update. The full gradient regularized scheme is detailed below in \cref{alg:boostingGradRegNewton} (the implicit exact step is denoted in parentheses):
\begin{algorithm}[H]
\caption{Gradient-Regularized Restricted Newton}
\label{alg:boostingGradRegNewton}
\begin{algorithmic}[1]
    \STATE \textbf{Input:} Loss $L:\mathcal{H}\to\bbR$ with $2M$-Lipschitz Hessian. Learning rate $\eta$. Weak learner family $\mathcal{F}$. Initial $F_0$.
    \FOR{$k = 0, 1, 2, \dots$}
        \STATE $g_k, \: H_k \leftarrow \nabla L(F_k), \: \nabla^2L(F_k)$
        \STATE Choose $\lambda_k \geq \sqrt{M\|g_k\|}$
        %\STATE $\big($Exact target $f_{k+1} = \argmin_{f \in \mathcal{H}} Q_k^{reg}(f) = - (H_k + \lambda_k I)^{-1}[g_k] \big)$
        \STATE $\big($Exact target $f_{k+1} \leftarrow - (H_k + \lambda_k I)^{-1}g_k \big)$
        \STATE Weak learner $f^{w}_{k+1} \leftarrow \argmin_{f \in \mathcal{F}} Q_k^{reg}(f)$
        \STATE $F_{k+1} \leftarrow F_k + \eta f^{w}_{k+1}$
    \ENDFOR
\end{algorithmic}
\end{algorithm}




\subsection{Properties of Gradient Regularization}
We begin by establishing some basic properties the gradient regularized scheme. In the subsequent analysis, define $K_k := H_k + \lambda_k I$.

\begin{lemma}\label{lemmaLambdafMinusGkHk}
    The exact regularized update \eqref{eqRegBoostingRoundExact} satisfies
    \begin{equation}\label{eqRegNewtonUpdateDirection}
        \lambda_k f_{k+1} = -\left( g_k + H_k[f_{k+1}] \right).
    \end{equation}
\end{lemma}
\begin{proof}
    Multiply the Newton update \eqref{eqRegBoostingRoundExact} by $(H_k + \lambda_k I)$.
\end{proof}



\begin{lemma}\label{lemmaUsefulProperties}
    If $L$ is convex, then for all $u\in \mathcal{H}$
    \begin{align*}
        \|u\|_{K_k}^2  
        &:= \langle(H_k + \lambda_k I)u, u\rangle =  \lambda_k \|u\|^2 + \langle u, H_k u\rangle \\
        &\geq \lambda_k \|u\|^2.
    \end{align*}
    Furthermore, if the weak learner family $\mathcal{F}$ is closed under scalar multiplication, then the iterates of \cref{alg:boostingGradRegNewton} satisfy
    \begin{enumerate}
        \item $ \|f_{k+1}^w\|_{K_k}^2 = -\langle g_k, f_{k+1}^w \rangle$,
        \item $\lambda_k \|f_{k+1}^w\| \leq \|g_k\|$.
    \end{enumerate}
\end{lemma}
\begin{proof}
    The first result follows by convexity. For the rest, let
    $$\phi(t) := Q_k^{reg}(t f_{k+1}^w) = \frac{1}{2}t^2 \|f_{k+1}^w\|_{K_k}^2 + t \langle g_k, f_{k+1}^w\rangle.$$
    Differentiating gives $0 = \phi'(t) = t \|f_{k+1}^w\|_{K_k}^2 + \langle g_k, f_{k+1}^w\rangle$. Since $\mathcal{F}$ was assumed to be scalable, the minimum must be attained at  $t=1$, and the result follows. As for the last inequality, this follows by the first inequality and Cauchy-Schwartz:
    \begin{align*}
        \lambda_k \|f_{k+1}^w\|^2 \leq \|f_{k+1}^w\|_{K_k}^2 = -\langle g_k, f_{k+1}^w \rangle \leq \|g_k\| \|f_{k+1}^w\|. 
    \end{align*}
\end{proof}

%\vspace{-10pt}
We impose a standard regularity condition on the loss $L$ in terms of the Lipschitz continuity of its Hessian, as is common for regularized Newton methods \cite{2006CubicRegularizationOfNewtonMethodAndItsGlobalPerformance, 2008AcceleratingTheCubicRegularizationOfNewtonsMethodOnConvexProblems, 2023RegularizedNewtonMethodWithGlobalConvergence, 2024GradientRegularizationOfNewtonMethodWithBregmanDistances}.

\begin{assumption}\label{assumptionHessLipAndCosine}
    The Fr\'echet Hessian $\nabla^2 L : \mathcal{H} \to \mathcal{L(H)}$ is $2M$-Lipschitz continuous for some constant $M>0$:
    \begin{align*}
        \|\nabla^2L(f)-\nabla^2L(g)\| \leq 2M\|f-g\| \qquad \forall f,g\in \mathcal{H}.
    \end{align*}
\end{assumption}

Before we delve into the convergence rate analysis, we recall the following standard results on Lipschitz Hessians, well known in the finite-dimensional $\mathbb{R}^d$ setting. See \cref{app:SmoothnessLemmas} for proof.

\begin{lemma}\label{lemmaStandardHessianResults}
    If $\nabla^2 L$ is $2M$-Lipschitz, then for all $f,g\in \mathcal{H}$
    \begin{enumerate}
        \item $L(f) \leq L(g) + \left\langle \nabla L(g), f-g\right\rangle + \frac{1}{2}\langle \nabla^2 L(g)(f-g), f-g\rangle + \frac{M}{3}\|f-g\|^3$,
        \item  $\|\nabla L(f) - \nabla L(g) - \nabla^2 L(g)(f-g)\| \leq M\|g-f\|^2$.
    \end{enumerate}
\end{lemma}

The following lemma shows that the weak iterates of \cref{alg:boostingGradRegNewton} guarantee a strict decrease in loss. From this point on, we assume that $\mathcal{F}$ is closed under scalar multiplication.

\begin{lemma}\label{lemmaDecreaseInLossNEW}
If $L$ is convex with $2M$-Lipschitz Hessian, and $\eta\leq 2$, then the iterates of \cref{alg:boostingGradRegNewton} satisfy 
\begin{align}
    L(F_k) - L(F_{k+1}) 
    %&\geq  \left( \eta - \frac{\eta^2}{2} - \frac{\eta^3}{3} \right) \|f_{k+1}^w\|_{K_k}^2 + \frac{\eta^2}{2}\lambda_k \|f^w_{k+1}\|^2 \\
    &\geq \left( \eta - \frac{\eta^3}{3} \right) \lambda_k \|f^w_{k+1}\|^2
\end{align}
\end{lemma}
\begin{proof}
    Using \cref{lemmaStandardHessianResults} and \cref{lemmaUsefulProperties}, we write
    \begin{align*}
        L&(F_{k+1}) - L(F_k) \\
        &\leq  \eta \langle  g_k,  f_{k+1}^w\rangle + \frac{\eta^2}{2}\langle H_k(f_{k+1}^w),  f_{k+1}^w\rangle + \frac{M}{3} \eta^3\|f_{k+1}^w\|^3 \\
        &= -\eta \|f_{k+1}^w\|_{K_k}^2 + \frac{\eta^2}{2}\|f_{k+1}^w\|_{K_k}^2 - \frac{\eta^2}{2}\lambda_k\|f_{k+1}^w\|^2 + \frac{M}{3}\eta^3 \|f_{k+1}^w\|^3 \\
        &= -\left( \eta - \frac{\eta^2}{2} \right) \|f_{k+1}^w\|_{K_k}^2 - \frac{\eta^2}{2}\lambda_k\|f_{k+1}^w\|^2 +  \frac{M}{3}\eta^3 \|f_{k+1}^w\|^3.
    \end{align*}
    By \cref{lemmaUsefulProperties} and $\lambda_k \geq \sqrt{M \|g_k\|}$ we see that
    \begin{align*}
        M \|f_{k+1}^w\|^3
        \leq \frac{M \|g_k\|}{\lambda_k}\|f_{k+1}^w\|^2
        \leq \lambda_k\|f_{k+1}^w\|^2.
    \end{align*}
    Combining the above, and again using \cref{lemmaUsefulProperties}, gives the result.
\end{proof}


\subsection{Convergence of Restricted GRN}

We now require a way to quantify how well the weak updates $f^w_{k+1}$ approximate the exact updates $f_{k+1}$ in the gradient regularized scheme. Previously in \cref{lemmaStandardSmoothnessResults} for vanilla Newton, the decrease in loss was determined by the Hessian-induced norm $\|\cdot\|_{H_k}$. In contrast, for the regularized scheme the decrease in loss is governed by the quantity $\lambda_k \|f^w_{k+1}\|$. As such, the cosine angle from \cref{defCosineAngle} is no longer well suited for theoretical analysis of GRN, since the scheme and its convergence proof are fundamentally driven by gradient growth rather than directional alignment. Consequently, motivated by \cref{lemmaLambdafMinusGkHk}, we introduce the notion of an implied weak gradient $g^w_k$, together with a corresponding weak gradient edge condition. This condition essentially states that the weak gradient remains within a relative radius of the true gradient, at each iteration $k$.

\begin{definition}\label{defWeakGradientEdge}
    We define the \textbf{implied weak gradient} as
    \begin{align}\label{eqWeakGradient}
        g_k^w : = - (H_k + \lambda_k I)f_{k+1}^w.
    \end{align}
    We say that the iterates $f^w_{k+1}$ have \textbf{weak gradient edge} $\gamma$ if there exists a real number $0<\gamma \leq 1$ satisfying 
    \begin{align}\label{eqWeakGradientEdge}
        \|g_k^w - g_k\|^2 \leq (1-\gamma^2) \|g_k\|^2.
    \end{align}
\end{definition}

The following lemma relates the gradient at consecutive iterations and serves as the primary tool for our convergence analysis.

\begin{lemma}[\textbf{Gradient Growth}]\label{lemmagTProperties}
    If the iterates of $\cref{alg:boostingGradRegNewton}$ have weak gradient edge $\gamma$, then
    \begin{align}\label{eqGradientGrowthMainPaper}
        \|g_{k+1}\|
        &\leq \left(1 + \eta^2 + \eta \sqrt{1-\gamma^2}\right)\|g_k\|.
    \end{align}
\end{lemma}
\begin{proof}[Proof Sketch]
    Motivated by \cref{lemmaStandardHessianResults}, we use the identity
    \begin{align}
        g_{k+1} 
        &= g_k + \eta H_k f^w_{k+1} + E \\
        &= g_k + \eta (-g_k^w-\lambda_k f_{k+1}^w) + E \\
        &= (1-\eta) g_k + \eta(g_k - g_k^w) - \eta \lambda_k f^w_{k+1} + E\label{eqGradientGrowth}
    \end{align}
    where $\|E\| \leq M\eta^2 \|f^w_{k+1}\|^2$. Applying the triangle inequality, the weak learner edge, and \cref{lemmaUsefulProperties} yields the result. In the upcoming convergence rate proofs, both the final bound \eqref{eqGradientGrowthMainPaper} and intermediate expressions using \eqref{eqGradientGrowth} play an important role, see \cref{app:GlobalConvergenceProofs} for proofs.
\end{proof}


We now present our main result: establishing the global convergence of gradient regularized Newton's method with weak learners. Despite the inexact updates inherent to boosting and Restricted Newton Descent, \cref{alg:boostingGradRegNewton} retains the $\mathcal{O}(\frac{1}{k^2})$ rate, matching first-order momentum boosting \cite{2020AcceleratingGradientBoostingMachines} without the requirement of fitting two trees per iteration. We use the following standard definition from the cubic Newton literature.

\begin{definition}
    $L$ has \textbf{finite sublevel sets} if a finite optimum $F^* \in \argmin_{F \in \mathcal{H}} L(F)$ exists, and if there exists a constant $\Lambda>0$ such that $L(F) \leq L(F_0) \implies \|F-F^*\| \leq \Lambda$.
\end{definition}

\begin{theorem}[\textbf{Global Rate II}]\label{thm:global_rate_II}
    Assume $L$ is convex with $2M$-Lipschitz Hessian, finite sublevel sets with diameter $\Lambda$, and weak gradient edge $\gamma$. Let $\lambda_k = C\sqrt{M\|g_k\|}$ for some $C\geq 1$. For any learning rate $\eta \in (0,1]$, the gradient regularized Newton scheme has global convergence rate
    \begin{align*}
        L(&F_k) - L(F^*) \leq \\
        &\mathcal{O}\max\left(\frac{C^2 M \Lambda^3}{\eta^2 k^2 \gamma^{12}},\:\: \Lambda\exp\left(-k\eta\gamma^2\big(\frac{\gamma^2}{2}+2-2\eta\big) / 8\right)\right).
    \end{align*}
\end{theorem}
\begin{proof}[Proof Sketch] We generalize the analysis of \citet{2023RegularizedNewtonMethodWithGlobalConvergence} to the Hilbert space setting with inexact weak learners. The proof proceeds by analyzing two regimes based on the fraction of iterates where the gradient fails to contract sufficiently, determined by $\eta$ and $\gamma$. In the first regime, we derive a recurrence of the form $\alpha_{k+1} \leq \alpha_k - \tau \alpha_k^{3/2}$ for some $\tau > 0$, yielding a $\mathcal{O}(\frac{1}{\eta^2 \gamma^{12} k^2})$ rate. In the second regime, the cumulative contraction dominates, resulting in a linear rate. The global rate is governed by the maximum of these two bounds; see \cref{app:GlobalConvergenceProofs} for more details.
\end{proof}





\begin{theorem}[\textbf{Local Rate}]\label{thm:local_rate}
Let $\lambda_k = C\sqrt{M\|g_k\|}$ for some $C\geq 1$. If $L$ is $\mu$-strongly convex with $2M$-Lipschitz Hessian, weak gradient edge $\gamma$, and if there exists a $k_0\geq 0$ such that $\|g_{k_0}\| \leq \frac{\mu^2}{4C^2M}\left( \frac{1-\sqrt{1-\gamma^2}}{1+\sqrt{1-\gamma^2}}\right)^2$, then for all $k\geq k_0$ we have
\vspace{-10pt}
\begin{align}
    \|g_{k+1}\| \leq \frac{3 + \rho}{4} \|g_k\|,
\end{align}
where $\rho = 1 - \eta(1-\sqrt{1-\gamma^2}) = \mathcal{O}(1-\eta\frac{ \gamma^2}{2})$. Hence the iterates of \cref{alg:boostingGradRegNewton} converge linearly locally.
\end{theorem}
\begin{proof}[Proof Sketch]\vspace{-3pt}
    The weak gradient edge implies that $\|g_k^w\| \leq \left( 1 + \sqrt{1-\gamma^2} \right) \|g_k\| := C_\gamma \|g_k\|$, while the $\mu$-strong convexity implies that
    \vspace{-2pt}
    \begin{align*}
        \|f^w_{k+1}\| 
        = \|(H_k+\lambda I)^{-1} g^w_k\|
        \leq \frac{1}{\mu}\|g^w\|
        \leq \frac{C_\gamma}{\mu}\|g^w\|.
    \end{align*}
    \vspace{-2pt}Combining these bounds with \cref{lemmaUsefulProperties,lemmagTProperties} yields the recurrence
    \vspace{-2pt}
    \begin{align*}
        \|g_{k+1}\|
        &\leq \rho \|g_k\| + \sqrt{C^2 M \eta^2 \frac{C_\gamma^2}{\mu^2}} \|g_k\|^\frac{3}{2} + C^2 M \eta^2 \frac{C_\gamma^2}{\mu^2}\|g_k\|^2.
    \end{align*}
    \vspace{-2pt}The result follows by solving for the region of strict contraction $\|g_{k+1}\| < \|g_k\|$, where the dependence on $\eta$ cancels out. See \cref{app:GlobalConvergenceProofs} for details. 
\end{proof}
\begin{remark}
    If $\eta=1$ and $\gamma=1$, the coefficient of the linear $\|g_k\|$ term vanishes, allowing for a superlinear rate with exponent $\frac{3}{2}$, recovering the classical result of \citet{2023RegularizedNewtonMethodWithGlobalConvergence}. However, when $\eta < 1$ or $\gamma < 1$, this term is non-zero, and we are only able to prove a linear rate. We conjecture that superlinear rates are generally unattainable in Newton schemes involving weak iterates, suggesting an interesting direction for future research.
\end{remark}


\section{Practical Considerations}

In this section, we discuss when the regularity assumptions of \cref{lemmaStandardSmoothnessResults,thm:global_rate_II,thm:local_rate} are satisfied in practice for boosting, and how the GRN scheme can be implemented into existing tree boosting algorithms with negligible computational overhead.

\subsection{Regularity of the Boosting Objective}

Our convergence analysis relies on standard convex optimization assumptions, such as Lipschitz-continuous Hessians, smoothness, and strong convexity. Using the same notation as \cref{subsec:GBDTsAsRestrictedConvexOptimization}, because the boosting empirical risk $L(F) = \frac{1}{N} \sum_{i=1}^N l(F(x_i), y_i)$ is constructed from a pre-specified element-wise loss $l$ (e.g., MSE, cross entropy, Charbonnier), one can prove that the regularity of $l$ carries over to $L$ in the boosting space $\mathcal{H} = \mathcal{L}^2(\widehat{\nu}_{X,N})$. We formalize this result below. Note, however, that one must be careful when dealing with the previously mentioned Lipschitz constants: the regularity conditions of $l$ are defined in the Euclidean geometry, while the Lipschitz constants used by \cref{lemmaStandardSmoothnessResults,thm:global_rate_II,thm:local_rate} are defined in the boosting space with its own intrinsic geometry and inner products. 

\vspace{5pt}
\begin{proposition}\label{prop:regularity}
    If the element-wise loss $l(\cdot, y_i)$ has an $M_0$-Lipschitz Hessian in $\mathbb{R}^K$ uniformly across all fixed targets $y_i$, then the empirical risk $L$ has an $M_0\sqrt{N}$-Lipschitz Hessian in the boosting space $\mathcal{L}^2(\widehat{\nu}_{X,N})$. Additionally, if $l$ is $S$-smooth or $\mu$-strongly convex, then $L$ is also $S$-smooth or $\mu$-strongly convex, respectively.
\end{proposition}

The proof can be found in \cref{app:RegularityAssumptionsCarryOverProof}. This result ensures that if the base loss $l$ is well-behaved, the boosting objective is as well. Moreover, because $l$ is chosen prior to training, its Lipschitz constant $M_0$ can be computed analytically. This avoids the need for adaptive backtracking line search algorithms to estimate this constant, as is typically required in finite-dimensional convex optimization \cite{2022ADampedNewtonMethodAchievesGlobalMinus2AndLocalQuadraticConvergenceRate,2023RegularizedNewtonMethodWithGlobalConvergence,2025OPTAMI}.

\subsection{Implementing GRN for GBDTs}
The GRN scheme is very simple to implement in practice and only differs from vanilla Newton boosting by an extra additive $\ell_2$-regularization term based on the gradient norm. Both methods use the exact same split finding, gain calculation, and leaf fitting formulas. The only difference lies in the effective $\ell_2$-regularization used. Standard tree boosting implementations (e.g., XGBoost, LightGBM, CatBoost) use a static $\ell_2$-regularization $\lambda_{\text{base}}$. GRN, on the other hand, uses an adaptive regularization term $\lambda_k$ calculated via
\begin{equation}
    \lambda_k = \lambda_{\text{base}} + \sqrt{M \|g_k\|_{\mathcal{H}}}.
\end{equation}
Using the notation of \cref{subsec:GBDTsAsRestrictedConvexOptimization}, the exact gradient norm in the boosting space is computed from the sample gradients $g_{(i)}$ as $\|g_k\|_{\mathcal{H}} = \big(\frac{1}{N}\sum_{i=1}^N \|g_{(i)}\|^2_{\mathbb{R}^K}\big)^{1/2}$. Because updating this adaptive $\ell_2$ term only requires a single vector norm calculation and a scalar addition, and the main computational bottleneck for GBDT algorithms is split finding, the computational overhead of GRN is negligible.


\section{Numerical Experiments}\label{secNumericalResults}
In this section, we empirically validate our theoretical findings regarding the convergence of boosting schemes and the behaviour of weak learner conditions. We implement first-order gradient boosting, Newton boosting, and the Gradient Regularized Newton (GRN) boosting scheme in JAX \cite{2018jaxGithub}, following the standard depth-wise greedy tree construction strategy used by XGBoost \cite{2016XGBoost}. The code is provided in the supplementary material and at \url{github.com/nikitazozoulenko/global-newton-boosting}.



\subsection{Charbonnier loss divergence and convergence}
The Charbonnier loss, defined as $l(\widehat{y}, y) = \sqrt{1+(y-\widehat{y})^2}$ $-1$, is a classical example of a loss function where Newton's method can diverge (see \cref{app:newton_divergence}). We investigate if this also is true for boosting, using the Wine Quality UCI dataset \cite{UCIWineQuality} with depth 4 trees. For GRN we set $M=1$, and otherwise use a roughly equivalent static $\ell_2$-regularization of $1.0$ when explicitly stated. \Cref{fig:Charbonnier} compares the training loss of Newton boosting, first-order gradient boosting, and GRN for different learning rates $\eta$. For $\eta=1$, vanilla Newton boosting diverges, whereas GRN decreases the fastest. For a smaller $\eta=0.1$, GRN initially tracks the trajectory of fixed $\ell_2$-regularized Newton boosting but eventually outperforms it once the gradient norm decreases sufficiently. Vanilla Newton boosting minimizes the loss most rapidly in this scenario, but is limited to smaller learning rates. Notably, $\ell_2$-regularized models will converge to a regularized optimum, whereas GRN provably minimizes the original non-regularized loss.


\begin{figure}[t]
  \vskip -0.05in
  \begin{center}
    \centerline{\includegraphics[width=\columnwidth]{Figures/charbonnier_train_loss_paper.pdf}}
    \caption{
      Charbonnier train loss on the Wine Quality UCI dataset, for different gradient boosting schemes and learning rates $\eta$. Vanilla Newton boosting diverges for $\eta=1$.
    }
    \label{fig:Charbonnier}
  \end{center}
  \vskip -0.3in
\end{figure}



\subsection{Weak learner conditions and tree depth}

To better understand the theoretical constants governing convergence, we train models on a subset of the Higgs dataset \cite{UCIHiggs} using BCE loss with $\ell_2$-regularization. We compute the cosine angle $\Theta_k$ and weak learner edge $\gamma_k$ in \cref{fig:bceTreeDepth} at every iteration $k$, via \eqref{eqHessianCosineAngle} and \eqref{eqWeakGradientEdge}, by instantiating $f^w_{k+1}, f_{k+1}, g^w_k$, and $g_k$ as vectors in $\mathbb{R}^N$. As expected, increasing tree depth results in stronger weak learners, reflected in higher values for $\Theta_k$ and $\gamma_k$. We observe that these metrics are high during the initial iterations of boosting and eventually plateau at positive values. The trajectories of the cosine angle $\Theta_k$ and weak gradient edge $\gamma_k$ are remarkably similar, which is somewhat unexpected given that they are defined in different geometries. Conceptually, $\gamma_k$ (\cref{defWeakGradientEdge}) can be viewed as an approximate $(H_k + \lambda_k I)^2$-induced cosine angle between $f^w_{k+1}$ and $f_{k+1}$, in contrast to the standard $H_k$-induced cosine angle $\Theta_k$. If $f^w_{k+1}$ and $f_{k+1}$ were orthogonal under the $(H_k + \lambda_k I)^2$-induced norm, $\gamma_k$ would constitute a true cosine angle; however, we cannot theoretically guarantee this at present.

A similar relationship to tree depth holds for dataset size; larger datasets generally imply harder optimization problems, resulting in smaller cosine angles for a fixed tree depth. Furthermore, as optimization progresses, fitting the diminishing residuals becomes increasingly difficult for the weak learners, naturally leading to lower cosine angles.



\begin{figure}[t]
  \vskip -0.0in
  \begin{center}
    \centerline{\includegraphics[width=\columnwidth]{Figures/higgs_metrics.pdf}}
    \caption{
      Binary cross entropy, cosine angle $\Theta_k$, and weak gradient edge $\gamma_k$ computed at each boosting iteration $k$, for varying tree depths $d$. The exact values are noisy between iterations, but the rolling window average (highlighted in bold) remains stable.
    }
    \label{fig:bceTreeDepth}
  \end{center}
  \vskip -0.0in
\end{figure}











\section{Conclusion and Future Work}

We introduced Restricted Newton Descent to analyze the global convergence of vanilla Newton boosting and a novel gradient regularized Newton boosting scheme, based on the cosine angle $\Theta$ and weak gradient edge $\gamma$. For vanilla Newton boosting, we proved a linear convergence rate for Hessian-dominated losses. For the gradient regularized scheme, we established a global $\mathcal{O}(1/k^2)$ rate for general convex losses with Lipschitz Hessians. Our results suggest several promising directions for future research. First, given that finite-dimensional cubic and gradient regularized Newton can be accelerated to $\mathcal{O}(1/k^3)$, investigating whether such acceleration can be realized in the boosting setting is a natural next step. Second, it would be interesting to study the Restriction operation for other optimization schemes to determine which convergence properties are retained under weak updates, and see if the dependence on $\Theta^2$ and $\gamma^{12}$ can be improved. Finally, investigating the possibility of superlinear rates with weak learners and establishing generalization bounds for Newton boosting remain open problems. Future work will also explore other applications of Restricted Newton Descent beyond the gradient boosting setting.











\section*{Acknowledgements}
NZ has been supported by the Roth Scholarship at Imperial College London. DF was supported by the Göran Gustafsson Foundation for Research in Natural Sciences and Medicine via Alexander Berglund's 2024 Göran Gustafsson Prize. The work of TC was supported in part by UK Research and Innovation (UKRI) through the Engineering and Physical Sciences Research Council (EPSRC) via Programme Grants [Grant No. UKRI1010: High order mathematical and computational infrastructure for streamed data that enhance contemporary generative and large language models]. We acknowledge computational resources and support provided by the Imperial College Research Computing Service (DOI: \texttt{10.14469/hpc/2232}). For the purpose of open access, the authors have applied a Creative Commons Attribution (CC BY) licence to any Author Accepted Manuscript version arising.





\bibliography{preprint_main}
\bibliographystyle{Styles/icml2026}


%%%%%%%%%%%%%%%%%%%%%%%%%%%%%%%%%%%%%%%%%%%%%%%%%%%%%%%%%%%%%%%%%%%%%%%%%%%%%%%
%%%%%%%%%%%%%%%%%%%%%%%%%%%%%%%%%%%%%%%%%%%%%%%%%%%%%%%%%%%%%%%%%%%%%%%%%%%%%%%
% APPENDIX
%%%%%%%%%%%%%%%%%%%%%%%%%%%%%%%%%%%%%%%%%%%%%%%%%%%%%%%%%%%%%%%%%%%%%%%%%%%%%%%
%%%%%%%%%%%%%%%%%%%%%%%%%%%%%%%%%%%%%%%%%%%%%%%%%%%%%%%%%%%%%%%%%%%%%%%%%%%%%%%
\newpage
\appendix
\onecolumn






\section{Discussions on Losses and Gradient Boosting Trees}

In this appendix, we provide explicit calculations for the gradients and Hessians of common loss functions used in boosting. We also provide the full derivation of the optimal weight update for decision trees within the Newton boosting framework, recovering the standard XGBoost formula.

\subsection{Common loss functions used in boosting}


We denote the prediction for the $i$-th data point as $u_i = F(x_i)$. The gradient and Hessian with respect to the prediction are denoted by $g_{(i)} \in \mathbb{R}^K$ and $h_{(i)} \in \mathbb{R}^{K \times K}$, respectively.

\textbf{Mean Squared Error (MSE):}
For the quadratic loss $l(\widehat{y}, y) = \frac{1}{2}(\widehat{y}-y)^2$, the gradient and Hessian are given by:
\begin{align*}
    g_{(i)} &= \widehat{y}_i - y_i, \\
    h_{(i)} &= 1.
\end{align*}
\textbf{Binary Cross Entropy (BCE):}
For binary classification, let $y_i \in \{0,1\}$ be the label and $u_i \in \mathbb{R}$ be the logit. The loss is defined as $l(u_i, y_i) = -[y_i \log \sigma(u_i) + (1-y_i) \log(1-\sigma(u_i))]$, where $\sigma(z) = (1+e^{-z})^{-1}$ is the sigmoid function. Using the chain rule, one finds that the derivatives are
\begin{align*}
    g_{(i)} &= \sigma(u_i) - y_i, \\
    h_{(i)} &= \sigma(u_i)(1-\sigma(u_i)).
\end{align*}
\textbf{Categorical Cross Entropy (CCE):}
In the multi-class setting ($K > 2$), let $y_i \in \{0,1\}^K$ be the one-hot encoded label and $u_i \in \mathbb{R}^K$ be the vector of logits. The loss is $l(u_i, y_i) = -\sum_{k=1}^K (y_i)_k \log (p_i)_k$, where $p_i = \operatorname{softmax}(u_i)$ is the vector of class probabilities. The gradient and Hessian are
\begin{align*}
    g_{(i)} &= p_i - y_i, \\
    h_{(i)} &= \text{diag}(p_i) - p_i p_i^\top,
\end{align*}
where $\text{diag}(p_i)$ is the diagonal matrix with elements of $p_i$ on the diagonal, and $p_i p_i^\top$ denotes the vector outer product.




\subsection{Newton Tree Boosting Derivation}\label{appendixXGBFormulaDerivation}


Let $f(x)$ be a decision tree defined by $f(x) = \sum_j w_j I_j(x)$, where $\{I_j\}_{j=1}^L$ is a disjoint partition of the feature space (i.e., the leaves of the tree) and $w_j \in \mathbb{R}^K$ are the leaf values. The second-order approximation of the loss, $Q(f)$, is given by
\begin{align*}
    Q(f) 
    &=\frac{1}{N} \sum_{i=1}^N \left\langle f(x_i), g_{(i)} + \frac{1}{2} h_{(i)} f(x_i) \right\rangle_{\mathbb{R}^K} \\
    &= \frac{1}{N} \sum_{i=1}^N \left\langle f(x_i), g_{(i)} \right\rangle_{\mathbb{R}^K} + \frac{1}{2N} \sum_{i=1}^N \left\langle f(x_i), h_{(i)} f(x_i) \right\rangle_{\mathbb{R}^K}.
\end{align*}
We analyze the linear and quadratic terms separately. Substituting the definition of the tree $f(x)$, we find that
\begin{align*}
    \sum_{i=1}^N \left\langle f(x_i), g_{(i)} \right\rangle_{\mathbb{R}^K} 
    &= \sum_{i=1}^N \left\langle \sum_j w_j I_j(x_i), g_{(i)} \right\rangle_{\mathbb{R}^K}  \\
    &= \sum_j \left\langle w_j, \sum_{i=1}^N I_j(x_i) g_{(i)}\right\rangle_{\mathbb{R}^K} \\
    &= \sum_j \left\langle w_j, \sum_{i \in I_j} g_{(i)}\right\rangle_{\mathbb{R}^K}.
\end{align*}
For the quadratic term, we use the property that the regions are disjoint, implying $I_j(x)I_k(x) = \delta_{jk}I_j(x)$, where $\delta_{jk}$ is the Kronecker delta. 
\begin{align*}
    \sum_{i=1}^N \left\langle f(x_i), h_{(i)} f(x_i) \right\rangle_{\mathbb{R}^K} 
    &= \sum_{i=1}^N \left\langle \sum_j w_j I_j(x_i), h_{(i)} \sum_k w_k I_k(x_i) \right\rangle_{\mathbb{R}^K} \\
    &= \sum_{i=1}^N \sum_j \sum_k I_j(x_i) I_k(x_i) \left\langle w_j, h_{(i)} w_k \right\rangle_{\mathbb{R}^K} \\
    &= \sum_{i=1}^N \sum_j I_j(x_i) \left\langle w_j, h_{(i)} w_j \right\rangle_{\mathbb{R}^K} \\
    &= \sum_j \sum_{i \in I_j} \left\langle w_j, h_{(i)} w_j \right\rangle_{\mathbb{R}^K} \\
    &= \sum_j \left\langle w_j, \left( \sum_{i \in I_j} h_{(i)} \right) w_j \right\rangle_{\mathbb{R}^K}.
\end{align*}
The total objective decomposes into a sum of independent objectives for each leaf $j$. Differentiating with respect to a specific weight $w_j$, we obtain the local quadratic
\begin{align*}
    q(w_j) =  \left\langle w_j, \sum_{i \in I_j} g_{(i)} \right\rangle 
    + \frac{1}{2} \left\langle w_j, \left( \sum_{i \in I_j} h_{(i)} \right) w_j \right\rangle.
\end{align*}
Taking the gradient with respect to $w_j$ yields
\begin{align*}
    \nabla q(w_j)
    &= \sum_{i \in I_j} g_{(i)} + \left( \sum_{i \in I_j} h_{(i)} \right) w_j.
\end{align*}
Setting the gradient to zero gives the optimal weight update
\begin{align*}
    w_j^* &=  -\left( \sum_{i \in I_j} h_{(i)} \right)^{-1} \sum_{i \in I_j} g_{(i)}.
\end{align*}
Substituting $w_j^*$ back into the objective (and defining $G_j = \sum_{i \in I_j} g_{(i)}$ and $H_j = \sum_{i \in I_j} h_{(i)}$), we recover the standard score used in Newton boosting and XGBoost:
\begin{align*}
    Q(f^*) = \sum_j q(w_j^*) 
    &=  -\frac{1}{2} \sum_j \left\langle H_j^{-1} G_j , G_j \right\rangle_{\mathbb{R}^K}.
\end{align*}



\begin{remark}
    The $\ell_2$-regularization used in standard implementations like XGBoost differs slightly from the one suggested by our theory. XGBoost adds a fixed scalar $\lambda$ to the aggregated leaf Hessian, yielding the update $w_j^* = -(H_j + \lambda)^{-1} G_j$. By contrast, our framework regularizes the weak learner in the Hilbert space $L^2(\widehat{\nu}_N)$. In this space, the regularization modifies the Hessian operator pointwise, which is equivalent to replacing each $h_{(i)}$ with $h_{(i)} + \lambda$. Summing these over the partition $I_j$ results in the update $w_j^* = -(H_j + |I_j| \lambda)^{-1} G_j$. Thus, the theoretically consistent regularization scales proportionally with the number of samples routed to the leaf. This is the formulation we use in our JAX implementation for the numerical experiments.
\end{remark}



% Let $f(x)$ be a decision tree defined by $f(x) = \sum_j w_j I_j(x)$, where $\{I_j\}_{j=1}^L$ is a disjoint partition of the feature space (i.e., the leaves of the tree) and $w_j \in \mathbb{R}^K$ are the leaf values. The second-order approximation of the loss, with explicit $\ell_2$ regularization $\lambda$, is given by
% \begin{align*}
%     Q^{reg}(f) 
%     &=\frac{1}{N} \sum_{i=1}^N \left\langle f(x_i), g_{(i)} + \frac{1}{2} h_{(i)} f(x_i) \right\rangle_{\mathbb{R}^K}
%     + \frac{\lambda}{2} \|f\|_H^2 \\
%     &= \frac{1}{N} \sum_{i=1}^N \left\langle f(x_i), g_{(i)} \right\rangle_{\mathbb{R}^K} + \frac{1}{2N} \sum_{i=1}^N \left\langle f(x_i), h_{(i)} f(x_i) \right\rangle_{\mathbb{R}^K}
%     + \frac{\lambda}{2N} \sum_{i=1}^N \|f(x_i)\|_{\mathbb{R}^K}^2.
% \end{align*}
% We analyze the linear, quadratic, and regularization terms separately. Substituting the definition of the tree $f(x)$, we find that
% \begin{align*}
%     \sum_{i=1}^N \left\langle f(x_i), g_{(i)} \right\rangle_{\mathbb{R}^K} 
%     &= \sum_{i=1}^N \left\langle \sum_j w_j I_j(x_i), g_{(i)} \right\rangle_{\mathbb{R}^K}  \\
%     &= \sum_j \left\langle w_j, \sum_{i=1}^N I_j(x_i) g_{(i)}\right\rangle_{\mathbb{R}^K} \\
%     &= \sum_j \left\langle w_j, \sum_{i \in I_j} g_{(i)}\right\rangle_{\mathbb{R}^K}.
% \end{align*}
% For the quadratic term, we use the property that the regions are disjoint, implying $I_j(x)I_k(x) = \delta_{jk}I_j(x)$, where $\delta_{jk}$ is the Kronecker delta. 
% \begin{align*}
%     \sum_{i=1}^N \left\langle f(x_i), h_{(i)} f(x_i) \right\rangle_{\mathbb{R}^K} 
%     &= \sum_{i=1}^N \left\langle \sum_j w_j I_j(x_i), h_{(i)} \sum_k w_k I_k(x_i) \right\rangle_{\mathbb{R}^K} \\
%     &= \sum_{i=1}^N \sum_j \sum_k I_j(x_i) I_k(x_i) \left\langle w_j, h_{(i)} w_k \right\rangle_{\mathbb{R}^K} \\
%     &= \sum_{i=1}^N \sum_j I_j(x_i) \left\langle w_j, h_{(i)} w_j \right\rangle_{\mathbb{R}^K} \\
%     &= \sum_j \sum_{i \in I_j} \left\langle w_j, h_{(i)} w_j \right\rangle_{\mathbb{R}^K} \\
%     &= \sum_j \left\langle w_j, \left( \sum_{i \in I_j} h_{(i)} \right) w_j \right\rangle_{\mathbb{R}^K}.
% \end{align*}
% For the regularization term, we similarly obtain
% \begin{align*}
%     \sum_{i=1}^N \|f(x_i)\|_{\mathbb{R}^K}^2
%     &= \sum_{i=1}^N \left\| \sum_j w_j I_j(x_i) \right\|_{\mathbb{R}^K}^2 \\
%     &= \sum_{i=1}^N \sum_j I_j(x_i) \|w_j\|_{\mathbb{R}^K}^2 \\
%     &= \sum_j \sum_{i \in I_j} \|w_j\|_{\mathbb{R}^K}^2 \\
%     &= \sum_j N_j \|w_j\|_{\mathbb{R}^K}^2,
% \end{align*}
% where $N_j = |I_j|$. The total objective decomposes into a sum of independent objectives for each leaf $j$. Differentiating with respect to a specific weight $w_j$, we obtain the local quadratic
% \begin{align*}
%     q(w_j) =  \left\langle w_j, \sum_{i \in I_j} g_{(i)} \right\rangle 
%     + \frac{1}{2} \left\langle w_j, \left( \sum_{i \in I_j} h_{(i)} \right) w_j \right\rangle
%     + \frac{1}{2} \lambda \frac{N_j}{N} \|w_j\|_{\mathbb{R}^K}^2.
% \end{align*}
% Taking the gradient with respect to $w_j$ yields
% \begin{align*}
%     \nabla q(w_j)
%     &= \sum_{i \in I_j} g_{(i)} + \left( \sum_{i \in I_j} h_{(i)} + \lambda \frac{N_j}{N} I \right) w_j.
% \end{align*}
% Setting the gradient to zero gives the optimal weight update
% \begin{align*}
%     w_j^* &=  -\left( \sum_{i \in I_j} h_{(i)} + \lambda \frac{N_j}{N} I \right)^{-1} \sum_{i \in I_j} g_{(i)}.
% \end{align*}
% Substituting $w_j^*$ back into the objective (and defining $G_j = \sum_{i \in I_j} g_{(i)}$ and $H_j = \sum_{i \in I_j} h_{(i)}$), we recover the regularized score
% \begin{align*}
%     Q(f^*) = \sum_j q(w_j^*) 
%     &=  -\frac{1}{2} \sum_j \left\langle \left(H_j + \lambda \frac{N_j}{N} I\right)^{-1} G_j , G_j \right\rangle_{\mathbb{R}^K}.
% \end{align*}

\clearpage



\section{Some Classical Optimization Results}\label{app:SmoothnessLemmas}

We collect several standard smoothness and Lipschitz-Hessian definitions and results from optimization theory \cite{2006CubicRegularizationOfNewtonMethodAndItsGlobalPerformance,2018LecturesOnConvexOptimization}, and present them in the Hilbert space setting.

\begin{definition}[$S$-smooth loss]
Let $\mathcal H$ be a real Hilbert space. A Fr\'echet differentiable function
$L : \mathcal H \to \mathbb R$ is called \emph{$S$-smooth} if its gradient is Lipschitz continuous
with constant $S$, i.e.,
\[
\|\nabla L(f) - \nabla L(g)\| \le S \|f - g\|
\quad \forall f,g \in \mathcal H.
\]
Equivalently, if $L$ is twice Fr\'echet differentiable, then its Hessian satisfies
\[
\|\nabla^2 L(h)\|_{\mathrm{op}} \le S
\quad \forall h \in \mathcal H,
\]
or equivalently $\nabla^2 L(h) \preceq S I$ as an operator inequality.
\end{definition}

\begin{definition}[$\mu$-strongly convex loss]
A function $L : \mathcal H \to \mathbb R$ is \emph{$\mu$-strongly convex} for some $\mu>0$ if one
of the following equivalent conditions holds:
\begin{enumerate}
    \item \textbf{(Quadratic lower bound):}
    The function $h \mapsto L(h) - \frac{\mu}{2}\|h\|^2$ is convex.
    \item \textbf{(First-order condition):}
    For all $f,g \in \mathcal H$,
    \[
    L(f) \ge L(g) + \langle \nabla L(g), f-g\rangle
    + \frac{\mu}{2}\|f-g\|^2.
    \]
    \item \textbf{(Gradient monotonicity):}
    For all $f,g \in \mathcal H$,
    \[
    \langle \nabla L(f) - \nabla L(g), f-g\rangle \ge \mu \|f-g\|^2.
    \]
    \item \textbf{(Second-order condition):}
    If $L$ is twice Fr\'echet differentiable, then
    \[
    \nabla^2 L(h) \succeq \mu I
    \quad \forall h \in \mathcal H.
    \]
\end{enumerate}
\end{definition}

\begin{definition}
The Fr\'echet Hessian $\nabla^2 L : \mathcal{H} \to \mathcal{L(H)}$ is $2M$-Lipschitz continuous if there exists a constant $M>0$ such that
    \begin{align*}
        \|\nabla^2L(f)-\nabla^2L(g)\| \leq 2M\|f-g\| \qquad \forall f,g\in \mathcal{H}.
    \end{align*}
\end{definition}


%TODO state in hilbert space setting? Change from L-smoothness to S-smoothness. function f to function L.

%\begin{definition}[$S$-smooth function]
%A differentiable function $f: \mathbb{R}^n \to \mathbb{R}$ is $S$-smooth if its gradient is Lipschitz continuous with constant $S$, i.e., $\|\nabla f(\mathbf{x}) - \nabla f(\mathbf{y})\| \le L \|\mathbf{x} - \mathbf{y}\|$ for all $\mathbf{x}, \mathbf{y} \in \mathbb{R}^n$. Equivalently, if $f$ is twice-differentiable, its Hessian is bounded, i.e., $\nabla^2 f(\mathbf{x}) \preceq S\mathbf{I}$.
%\end{definition}

%\begin{definition}[$\mu$-strongly Convex Function]
%A function $f: \mathbb{R}^n \to \mathbb{R}$ is said to be $\mu$-strongly convex with modulus $\mu > 0$ if one of the following equivalent conditions holds:

%\begin{enumerate}
%    \item \textbf{Quadratic Lower Bound:} The function $g(\mathbf{x}) = f(\mathbf{x}) - \frac{\mu}{2}\|\mathbf{x}\|^2$ is convex.
%    \item \textbf{First-Order Condition (for differentiable $f$):} For all $\mathbf{x}, \mathbf{y} \in \mathbb{R}^n$, the following inequality holds:
%    $$f(\mathbf{y}) \ge f(\mathbf{x}) + \nabla f(\mathbf{x})^T(\mathbf{y} - \mathbf{x}) + \frac{\mu}{2}\|\mathbf{y} - \mathbf{x}\|^2$$
%    \item \textbf{Monotonicity of Gradient (for differentiable $f$):} For all $\mathbf{x}, \mathbf{y} \in \mathbb{R}^n$,
%    $$(\nabla f(\mathbf{x}) - \nabla f(\mathbf{y}))^T(\mathbf{x} - \mathbf{y}) \ge \mu\|\mathbf{x} - \mathbf{y}\|^2$$
 %   \item \textbf{Second-Order Condition (for twice-differentiable $f$):} The Hessian matrix is bounded below by $\mu\mathbf{I}$, i.e., for all $\mathbf{x} \in \mathbb{R}^n$,
 %   $$\nabla^2 f(\mathbf{x}) \succeq \mu\mathbf{I}$$
 %   This is equivalent to stating that all eigenvalues of the Hessian satisfy %$\lambda_{\min}(\nabla^2 f(\mathbf{x})) \ge \mu$.
%\end{enumerate}
%\end{definition}



% In the classical convex optimization literature, Nesterov's momentum achieves a $O(1/k^2)$ global rate for $\sigma$-smooth convex functions. Newton's method with learning rate $\mu/\sigma$ has a global linear rate, assuming $\sigma$-smoothness and $\mu$-strongly convex. Without strong convexity, Newton's method doesnt converge globally. However, with cubic or gradient regularization, Newton's method is able to achieve a $O(1/k^2)$ global rate, assuming only $2M$-lipchitzness of the Hessian \cite{2006CubicRegularizationOfNewtonMethodAndItsGlobalPerformance, 2023RegularizedNewtonMethodWithGlobalConvergence, 2024GradientRegularizationOfNewtonMethodWithBregmanDistances}. This can be accelerated to $O(1/k^3)$ using the estimation sequences framework for cubic \cite{2008AcceleratingTheCubicRegularizationOfNewtonsMethodOnConvexProblems} and gradient regularization \cite{2024GradientRegularizationOfNewtonMethodWithBregmanDistances}


% \subsection{Rates for Boosting}

% Modern boosting papers use some kind of "minimal cosine angle" assumption between arbitrary vectors and the output of the best-fit weak learner (the definitions varies slightly between \cite{2020AcceleratingGradientBoostingMachines, 2020RandomizedGradientBoostingMachine, 2020SnapBoostAHeterogeneousBoostingMachine})
% \begin{definition}
%     The Minimal Cosine Angle (MCA) $\Theta$ is defined as the similarity between arbitrary vectors $r$ and the output of the best-fit weak learner $b_\tau(X)$
%     \begin{equation*}
%         \min_{r \in \bbR^n} \max_{\tau \in \calT} \operatorname{cosangle}{(r, b_\tau(X))}.
%     \end{equation*}
% \end{definition}
% It remains unclear whether one needs to view this via coordinate descent or via $L^2(P)$ inner products. Also, this definition potentially seems not quite right to me: the tree outputs $b_\tau(X)$ are limited by the range of the observed features. Can the inf be 0 since we are looking over all of $\bbR^n$, and not just plausible gradient/Newton targets.

% The boosting paper \cite{2020SnapBoostAHeterogeneousBoostingMachine} assumes that the loss $l$ is $S$-smooth and $\mu$-strongly convex, and obtain a linear rate of convergence, same as in \cite{2018GlobalLinearConvergenceOfNewtonsMethodWithoutStrongConvexityOrLipschitzGradients} but with extra $\Theta$ factors.

% The Google gradient boosting with momentum paper \cite{2020AcceleratingGradientBoostingMachines} obtains $O(1/T^2)$ rate assuming only $\sigma$-smoothness (and convexity? unclear). With $\mu$-strong convexity, they obtain a linear rate similar to the accelerated nesterov case in the classical literature.



% %Consider the binary cross entropy loss function for a single data point, $l(x) = -y \log(\sigma(x)) - (1-y) \log(1-\sigma(x))$, where $\sigma(x) = (1+e^{-x})^{-1}$ and $y \in \{0,1\}$.

% %The second derivative is given by $l''(x) = \sigma(x)(1-\sigma(x))$.

% %The third derivative is $l'''(x) = \sigma'(x)(1-2\sigma(x)) = \sigma(x)(1-\sigma(x))(1-2\sigma(x))$. The magnitude of this derivative is bounded, i.e., $|l'''(x)| \le \frac{1}{4}$. This implies that the Hessian is Lipschitz continuous.

% %However, as $|x| \to \infty$, we have $\sigma(x) \to 0$ or $\sigma(x) \to 1$, which means $l''(x) = \sigma(x)(1-\sigma(x)) \to 0$. Since the second derivative is not bounded below by a positive constant, the function is not strongly convex.




 

\begin{lemma}
    If the loss $L$ is $S$-smooth, then
    \begin{align}
        L(F_k+f) \leq L(F_k) + \langle g_k, f\rangle + \frac{S}{2} \|f\|^2.
    \end{align}
\end{lemma}
\begin{proof}
    Using Taylor's theorem with integral remainder, we express the expansion of the loss as
    \begin{equation}
        L(F_k+f) = L(F_k) + \langle g_k, f\rangle + \int_0^1 (1-t) \langle f, \nabla^2 L(F_k + tf) f \rangle \, dt.
    \end{equation}
    The definition of $S$-smoothness implies that the Hessian is bounded in operator norm by $S$, i.e., $\nabla^2 L(\cdot) \preceq S I$. Consequently, the quadratic form is bounded by $\langle f, \nabla^2 L(\cdot) f \rangle \leq S \|f\|^2$. Substituting this into the integral yields $\int_0^1 (1-t) S \|f\|^2 \, dt = \frac{S}{2}\|f\|^2$, completing the proof.
\end{proof}


\begin{lemma}
    If $\nabla^2 L$ is $2M$-Lipschitz, then for all $f,g\in \mathcal{H}$
    \begin{enumerate}
        \item $L(f) \leq L(g) + \langle \nabla L(g), f-g\rangle + \frac{1}{2}\langle \nabla^2 L(g)(f-g), f-g\rangle + \frac{M}{3}\|f-g\|^3$
        \item  $\|\nabla L(f) - \nabla L(g) - \nabla^2 L(g)(f-g)\| \leq M\|g-f\|^2$.
    \end{enumerate}
\end{lemma}
\begin{proof}
Let $h = f-g$. We begin with the Fundamental Theorem of Calculus and expand the terms:
\begin{align*}
L(f) - L(g) - \langle \nabla L(g), h \rangle
&= \int_0^1 \langle \nabla L(g+th) - \nabla L(g), h \rangle \,dt \\
&= \int_0^1 \left\langle \int_0^t \nabla^2 L(g+sh)[h]\,ds, h \right\rangle \,dt \\
&= \int_0^1 \int_0^t \langle h, \nabla^2 L(g+sh)[h] \rangle \,ds\,dt \\
&= \int_0^1 \int_0^t \langle h, \nabla^2 L(g)[h] \rangle \,ds\,dt + \int_0^1 \int_0^t \langle h, (\nabla^2 L(g+sh) - \nabla^2 L(g))[h] \rangle \,ds\,dt \\
&= \frac{1}{2} \langle h, \nabla^2 L(g)[h] \rangle + \int_0^1 \int_0^t \langle h, (\nabla^2 L(g+sh) - \nabla^2 L(g))[h] \rangle \,ds\,dt \\
&\leq \frac{1}{2} \langle h, \nabla^2 L(g)[h] \rangle + \int_0^1 \int_0^t \|h\| \cdot \|(\nabla^2 L(g+sh) - \nabla^2 L(g))[h]\| \,ds\,dt \\
&\leq \frac{1}{2} \langle h, \nabla^2 L(g)[h] \rangle + \int_0^1 \int_0^t \|\nabla^2 L(g+sh) - \nabla^2 L(g)\|_{op} \|h\|^2 \,ds\,dt \\
&\leq \frac{1}{2} \langle h, \nabla^2 L(g)[h] \rangle + \int_0^1 \int_0^t (2M s \|h\|) \|h\|^2 \,ds\,dt \\
&= \frac{1}{2} \langle h, \nabla^2 L(g)[h] \rangle + 2M\|h\|^3 \int_0^1 \frac{t^2}{2} \,dt \\
&= \frac{1}{2} \langle h, \nabla^2 L(g)[h] \rangle + \frac{M}{3}\|h\|^3.
\end{align*}
Rearranging the terms yields the desired first inequality. For inequality 2, we express the gradient difference as an integral of the Hessian.
\begin{align*}
    \|\nabla L(f) - \nabla L(g) - \nabla^2 L(g)[h]\|
    &= \left\| \int_0^1 \nabla^2 L(g+th)[h] \,dt - \int_0^1 \nabla^2 L(g)[h] \,dt \right\| \\
    &= \left\| \int_0^1 (\nabla^2 L(g+th) - \nabla^2 L(g))[h] \,dt \right\| \\
    &\leq \int_0^1 \| (\nabla^2 L(g+th) - \nabla^2 L(g))[h] \| \,dt \\
    &\leq \int_0^1 \| \nabla^2 L(g+th) - \nabla^2 L(g) \|_{op} \|h\| \,dt \\
    &\leq \int_0^1 (2M t \|h\|) \|h\| \,dt \\
    &= 2M \|h\|^2 \int_0^1 t \,dt \\
    &= M \|h\|^2.
\end{align*}
This completes the proof.
\end{proof}








\clearpage
\section{Proof of Regularity Conditions in the Boosting Space}
\label{app:RegularityAssumptionsCarryOverProof}

Recall from \cref{subsec:GBDTsAsRestrictedConvexOptimization} that we identify the empirical boosting space $\mathcal{H} = \mathcal{L}^2(\widehat{\nu}_{X,N})$ with $\mathbb{R}^{NK}$, equipped with the inner product $\langle f, h\rangle = \frac{1}{N}\sum_{i=1}^N \langle f(x_i),h(x_i) \rangle_{\mathbb{R}^K}$. The induced norm is therefore $\|f\| = \big(\frac{1}{N}\sum_{i=1}^N \|f(x_i)\|^2_{\mathbb{R}^K}\big)^{1/2}$.

\vspace{5pt}
\begin{proof}[Proof of \cref{prop:regularity}]
First, we prove the claim regarding the Lipschitz continuity of the Hessian. Suppose the base loss $l(u, y)$ has a $M_0$-Lipschitz Hessian with respect to the prediction $u \in \mathbb{R}^K$, meaning $\|\nabla_{11}^2 l(u, y) - \nabla_{11}^2 l(v, y)\|_{\text{op}} \leq M_0 \|u - v\|_{\mathbb{R}^K}$. In the gradient boosting space, we minimize the empirical risk $L(F) = \frac{1}{N} \sum_{i=1}^N l(F(x_i), y_i)$. Because the Hessian operator $\nabla^2 L(F)$ acts pointwise (i.e., it is block-diagonal over the $N$ samples), its operator norm on $\mathcal{H}$ is exactly the maximum of the operator norms of the individual $K \times K$ blocks. Thus, for any $F, G \in \mathcal{H}$, we can bound the difference in Hessians as follows:
\begin{align*}
    \|\nabla^2 L(F) - \nabla^2 L(G)\|_{\text{op}, \mathcal{H}} 
    &= \max_{1 \leq i \leq N} \|\nabla_{11}^2 l(F(x_i), y_i) - \nabla_{11}^2 l(G(x_i), y_i)\|_{\text{op}, \mathbb{R}^K} \\
    &\leq \max_{1 \leq i \leq N} M_0 \|F(x_i) - G(x_i)\|_{\mathbb{R}^K} \\
    &\leq M_0 \left( \sum_{i=1}^N \|F(x_i) - G(x_i)\|_{\mathbb{R}^K}^2 \right)^{1/2} \\
    &= M_0 \sqrt{N} \left( \frac{1}{N}\sum_{i=1}^N \|F(x_i) - G(x_i)\|_{\mathbb{R}^K}^2 \right)^{1/2} \\
    &= M_0 \sqrt{N} \|F - G\|.
\end{align*}
This proves that $L$ has an $M_0\sqrt{N}$-Lipschitz Hessian in the space $\mathcal{H}$.

Next, we prove that smoothness and strong convexity are preserved with the same constants. The Hessian quadratic form in the boosting space is given by $\langle f, \nabla^2 L(F)[f] \rangle = \frac{1}{N} \sum_{i=1}^N \langle f(x_i), \nabla_{11}^2 l(F(x_i), y_i) f(x_i) \rangle_{\mathbb{R}^K}$. If the base loss $l$ is $\mu$-strongly convex and $S$-smooth, its Euclidean quadratic form is bounded pointwise:
\begin{equation*}
    \mu \|f(x_i)\|_{\mathbb{R}^K}^2 \leq \langle f(x_i), \nabla_{11}^2 l(F(x_i), y_i) f(x_i) \rangle_{\mathbb{R}^K} \leq S \|f(x_i)\|_{\mathbb{R}^K}^2.
\end{equation*}
Summing these bounds over $i=1, \dots, N$ and dividing by $N$ directly yields $\mu \|f\|^2 \leq \langle f, \nabla^2 L(F)[f] \rangle \leq S \|f\|^2$.
\end{proof}









\clearpage
\section{Convergence Rate Proofs}\label{app:GlobalConvergenceProofs}

For ease of exposition, we restate the statement and result of the relevant lemmas and theorems from the main part of the paper.

\begin{lemma}\label{appendixLemmaGradientGrowth}
    If the iterates of \cref{alg:boostingGradRegNewton} have weak gradient edge $\gamma$, then the gradient growth is bounded by
    \begin{align}
        \|g_{k+1}\| 
        &= \left(1-\eta(1-\sqrt{1-\gamma^2})\right)\|g_k\| + (\eta+\eta^2)\lambda_k \|f_{k+1}^w\| \\
        &\leq \left(1 + \eta^2 + \eta \sqrt{1-\gamma^2}\right)\|g_k\|.
    \end{align}
\end{lemma}
\begin{proof}
    Using \cref{eqWeakGradient}, and \cref{lemmaLambdafMinusGkHk,lemmaStandardHessianResults}, we can write
    \begin{align}
        g_{k+1} 
        &= g_k + \eta H_k f^w_{k+1} + E \\
        &= g_k + \eta (-g_k^w-\lambda_k f_{k+1}^w) + E \\
        &= (1-\eta) g_k + \eta(g_k - g_k^w) - \eta \lambda_k f^w_{k+1} + E
    \end{align}
    where $\|E\| \leq M\eta^2 \|f^w_{k+1}\|^2$. Hence using the bounds in \cref{lemmaUsefulProperties} we obtain that
    \begin{align}
        \|g_{k+1}\| 
        &= \| (1-\eta) g_k + \eta(g_k - g_k^w) - \eta\lambda_k f^w_{k+1} + E \| \\
        &\leq (1-\eta) \|g_k\| + \eta \|g_k - g_k^w\| +  \eta \lambda_k \|f^w_{k+1}\| + M \eta^2 \|f^w_{k+1}\|^2 \label{eqGradientRecursionForLater}\\
        &\leq (1-\eta) \|g_k\| + \eta \sqrt{1-\gamma^2}\|g_k\| +  \eta \lambda_k \|f^w_{k+1}\| + \eta^2 \lambda_k \|f^w_{k+1}\| \\
        &= \left(1-\eta(1-\sqrt{1-\gamma^2})\right)\|g_k\| + (\eta+\eta^2)\lambda_k \|f_{k+1}^w\| \\
        &\leq \left(1 + \eta^2 + \eta \sqrt{1-\gamma^2}\right)\|g_k\|.
    \end{align}
\end{proof}







The following global convergence proof of the gradient regularized Newton scheme generalizes the proof of \citet{2023RegularizedNewtonMethodWithGlobalConvergence} to the setting of Restricted Newton Descent with a constant learning rate $\eta \in (0, 1]$ and weak learners. The primary challenge lies in handling the additional error terms introduced by the learning rate and the inexact weak iterates $f^w_{k+1}$, which we handled in our previous Lemmas in this paper. Additonally, the original proof which uses a fixed subsequence condition and fraction, we must make this dependent on both the learning rate $\eta$ and weak gradient edge $\gamma$.

\begin{theorem}\label{theoremAppendixGradRegGlobalConv}
    Assume $L$ is convex with $2M$-Lipschitz Hessian, finite sublevel sets with diameter $\Lambda$, and weak gradient edge $\gamma$. Let the regularization be chosen as $\lambda_k = C \sqrt{M\|g_k\|}$ for some $C\geq 1$. For any learning rate $\eta \in (0,1]$, the gradient regularized Newton scheme has the global convergence rate
   \begin{align}
        L(F_k) - L(F^*) \leq \max\left( 
            \underbrace{\Lambda \|g_0\| \exp\left( - \frac{k\eta}{8}(1-\sqrt{1-\gamma^2})(3-\sqrt{1-\gamma^2}-2\eta) \right)}_{\text{Linear Regime}}, \quad 
            \underbrace{\frac{1024 C^2 M \Lambda^3}{\eta^2 (1-\sqrt{1-\gamma^2})^6(2+k)^2}\frac{(1+\eta)^4}{(1-\frac{1}{3}\eta^2)^2}.}_{\mathcal{O}\left(\frac{1}{k^2}\right)\text{ Regime}}
        \right).
    \end{align}
\end{theorem}
\begin{proof}
    By \cref{lemmaDecreaseInLossNEW} we have that $L(F_k) \leq L(F_{k-1}) \leq ... \leq L(F_0)$, implying that the iterates $F_k$ remain within the sublevel set defined by $F_0$. Consequently, $\|F_k-F^*\| \leq \Lambda$ for all $T$. By convexity we have the bound
    \begin{align}
        L(F_k) - L(F^*)
        &\leq \langle g_k, F_k-F^*\rangle  \leq \|g_k\| \Lambda.
    \end{align}
    Let $\mathcal{I}_\infty := \{ t\in \mathbb{N} : \|g_{t+1}\|\geq (1-\eta a) \|g_t\| \}$ be the set of indices where the gradient does not contract sufficiently, for some $0<a<1$ to be specified later. Consider first any $k\in \mathcal{I}_\infty$. It follows from \cref{appendixLemmaGradientGrowth} that
    \begin{align}
        (1-\eta a) \|g_k\| \leq \|g_{k+1}\| \leq  \left(1-\eta(1-\sqrt{1-\gamma^2})\right)\|g_k\| + (\eta+\eta^2)\lambda_k \|f_{k+1}^w\|
    \end{align}
    Rearranging the terms we obtain that
    \begin{align}\label{eqCase1LambdafNormGradientBound}
        \lambda_k \|f_{k+1}^w\| \geq \frac{1 - \sqrt{1-\gamma^2} - a}{1+\eta} \|g_k\|
    \end{align}
    For the next argument we require the bound \eqref{eqCase1LambdafNormGradientBound} to be positive, hence we choose 
    \begin{align}
        a := \frac{1-\sqrt{1-\gamma^2}}{2}.
    \end{align}
    By \eqref{eqCase1LambdafNormGradientBound} and \cref{lemmaDecreaseInLossNEW} we have
    \begin{align}
        L(F_k) - L(F_{k+1})
        &\geq \left( 1 - \frac{1}{3}\eta^2 \right) \eta  \lambda_k \|f_{k+1}^w\|^2 \\
        &\geq \left( 1 - \frac{1}{3}\eta^2 \right) a^2 \frac{\eta}{(1+\eta)^2} \frac{\|g_k\|^2}{\lambda_k} \\
        &= \left( 1 - \frac{1}{3}\eta^2 \right) a^2  \frac{\eta}{(1+\eta)^2} \frac{\|g_k\|^\frac{3}{2}}{C\sqrt{M}} \\
        &\geq \left( 1 - \frac{1}{3}\eta^2 \right) \frac{a^2 }{C\sqrt{M\Lambda^3}} \frac{\eta}{(1+\eta)^2} \bigg(  L(F_k) - L(F^*) \bigg)^\frac{3}{2} \\
        &= \tau \bigg(  L(F_k) - L(F^*) \bigg)^\frac{3}{2}
    \end{align}
    where $\tau = \left( 1 - \frac{1}{3}\eta^2 \right) \frac{a^2 }{C\sqrt{M\Lambda^3}} \frac{\eta}{(1+\eta)^2}$. If this recursion held for all $k\geq 1$, the $\mathcal{O}(\frac{1}{k^2})$ rate would follow by \cref{lemmaAlpha3div2Recursion}. Instead, we consider the subsequence technique as in \cite{2023RegularizedNewtonMethodWithGlobalConvergence}. We start by enumerating $\mathcal{I}_\infty = \{t_0, t_1, t_2, ...\}$. We have that
    \begin{align}
        \alpha_{k+1} 
        &:=  L(F_{t_{k+1}}) - L(F^*) \\
        &= L(F_{t_k}) - L(F^*) +  L(F_{t_k+1}) - L(F_{t_k}) \\
        &\leq  L(F_{t_k}) - L(F^*) + \tau \bigg(L(F_{t_k}) - L(F^*) \bigg)^\frac{3}{2} \\
        &= \alpha_k - \tau \alpha_k^\frac{3}{2}
    \end{align}
    \cref{lemmaAlpha3div2Recursion} implies that $\alpha_{k} \leq \frac{4}{\tau^2(2+k)^2}$. Consider now any general $k\geq 1$. Define $\mathcal{I}_k = \{ t \in \mathcal{I}_\infty : t \leq k\}$. We consider two cases based on the cardinality of the index set $\mathcal{I}_k$, that is, $|\mathcal{I}_k| \geq pk$ and $|\mathcal{I}_k| < pk$, for some fraction $0<p<1$ to be specified later in the proof. 
    
    \textbf{Case} $|\mathcal{I}_k| \geq pk$: Using the fact that both $L(F_k)$ and $\alpha_k$ are decreasing sequences, we obtain that
    \begin{align}
        L(F_k) - L(F^*) 
        \leq L(F_{|\mathcal{I}_k|}) - L(F^*) 
        = \alpha_{|\mathcal{I}_k|} 
        \leq \alpha_{pk}
        \leq \frac{4}{\tau^2p^2(2+k)^2}
        = \frac{4 C^2 M \Lambda^3}{\eta^2 a^4p^2(2+k)^2}\frac{(1+\eta)^4}{(1-\frac{1}{3}\eta^2)^2}.
    \end{align}
    Consequently in this case we have a global $\mathcal{O}(\frac{1}{k^2})$ rate.
    
    \textbf{Case} $|\mathcal{I}_k| < pk$:
    For indices $t\notin \mathcal{I}_k$, we have the contraction $\|g_{t+1}\| \leq (1-\eta a) \|g_t\|$. For $t \in \mathcal{I}_k$ the gradient may grow, but it is bounded by \cref{appendixLemmaGradientGrowth} via $\|g_{t+1}\| \leq \left(1 + \eta^2 + \eta \sqrt{1-\gamma^2}\right) \|g_t\|$. Combining both we obtain that
    \begin{align}
        \|g_k\| 
        &\leq \left(1 + \eta^2 + \eta \sqrt{1-\gamma^2}\right)^{|\mathcal{I}_k|} (1-\eta a)^{k-|\mathcal{I}_k|} \|g_0\|\\
        &\leq \left(1 + \eta^2 + \eta \sqrt{1-\gamma^2}\right)^{pk}(1-\eta a)^{(1-p)k} \|g_0\|.
    \end{align}
    Using the bound $1+x \leq e^x$ and taking logarithms, we see that this scheme produces a contraction if
    \begin{align}
        p\left( \eta^2 + \eta \sqrt{1-\gamma^2} \right) - (1-p)\eta a < 0.
    \end{align}
    Rearranging the above, we find that we require
    \begin{align}
        p < \frac{a}{\eta + a + \sqrt{1-\gamma^2}}.
    \end{align}
    Substituting our choice $a = (1-\sqrt{1-\gamma^2})/2$, the denominator simplifies to $\eta + (1+\sqrt{1-\gamma^2})/2$. Since $\sqrt{1-\gamma^2} \leq 1$, this denominator is strictly less than $\eta + 1$. Therefore, for any $\eta \in (0, 1]$, the choice $p = \frac{a}{2}$ is sufficient. In general, one can show that $a$ and $p$ need to satisfy the inequality
    \begin{align}
        \frac{p\left(\eta+\sqrt{1-\gamma^2}\right)}{1-p} < a <1-\sqrt{1-\gamma^2}.
    \end{align}
    for the proof in both cases above to be valid. With the choice $4p = 2a = 1-\sqrt{1-\gamma^2}$, one finds through elementary algebra and the bound $1+x \leq e^x$ that the gradient growth can be bounded by
    \begin{align}
        \|g_k\| \leq \exp\left( - \frac{k\eta}{8}(1-\sqrt{1-\gamma^2})(3-\sqrt{1-\gamma^2}-2\eta) \right) \|g_0\|,
    \end{align}
    yielding a linear rate of convergence.
    
    \textbf{Final bound:} The final convergence rate is obtained by taking the maximum bound of the two cases above. Note that $1-\sqrt{1-\gamma^2} = \frac{\gamma^2}{2} + o(\gamma^2)$. Hence the global convergence rate in the first case is 
    \begin{equation}
        \mathcal{O}\left(\frac{C^2M \Lambda^3}{\eta^2 k^2 \gamma^{12}}\right),
    \end{equation}
    while for the second case this is 
    \begin{equation}
        \mathcal{O}\left(\exp\left(-k\eta\gamma^2\left(\frac{\gamma^2}{2} +2-2\eta\right) / 8\right)\right).
    \end{equation}
    This concludes the proof.
\end{proof}











The following lemma is a standard result in the convex optimization literature, see e.g. \cite{2006CubicRegularizationOfNewtonMethodAndItsGlobalPerformance}. We include the general proof here for completeness

\begin{lemma}\label{lemmaAlpha3div2Recursion}
    If $\alpha_k$ is a decreasing sequence satisfying $0 \leq \alpha_{k+1} \leq \alpha_k - c\alpha_k^{\frac{3}{2}}$ for some $c>0$, then for all $k \geq 0$:
    \begin{equation}
        \alpha_k \leq \frac{1}{\left( \frac{1}{\sqrt{\alpha_0}} + \frac{c}{2}k \right)^2} \leq \frac{4}{c^2 (k + 2)^2}.
    \end{equation}
\end{lemma}
\begin{proof}
    Using the hypothesis $\alpha_k - \alpha_{k+1} \geq c \alpha_k^{3/2}$, consider the increment
    \begin{align}
        \frac{1}{\sqrt{\alpha_{k+1}}} - \frac{1}{\sqrt{\alpha_k}}
        &= \frac{\sqrt{\alpha_k} - \sqrt{\alpha_{k+1}}}{\sqrt{\alpha_k \alpha_{k+1}}} \\
        &= \frac{\alpha_k - \alpha_{k+1}}{\sqrt{\alpha_k \alpha_{k+1}} (\sqrt{\alpha_k} + \sqrt{\alpha_{k+1}})} \\
        &\geq \frac{c \alpha_k^{3/2}}{\sqrt{\alpha_k \alpha_{k+1}} (\sqrt{\alpha_k} + \sqrt{\alpha_{k+1}})}.
    \end{align}
    Since the sequence is decreasing ($\alpha_{k+1} \leq \alpha_k$), we have $\sqrt{\alpha_{k+1}} \leq \sqrt{\alpha_k}$, hence
    \begin{align}
        \frac{1}{\sqrt{\alpha_{k+1}}} - \frac{1}{\sqrt{\alpha_k}}
        \geq \frac{c \alpha_k^{3/2}}{\sqrt{\alpha_k^2} (\sqrt{\alpha_k} + \sqrt{\alpha_k})}
        = \frac{c}{2}.
    \end{align}
    Summing from $i=0$ to $k-1$ yields a telescoping sum, from which we obtain that
    \begin{equation}
        \frac{1}{\sqrt{\alpha_k}} \geq \frac{1}{\sqrt{\alpha_0}} + \frac{ck}{2}.
    \end{equation}
    Inverting and squaring both sides yields the inequality
    \begin{equation}
        \alpha_k \leq \frac{1}{\left( \frac{1}{\sqrt{\alpha_0}} + \frac{c}{2}k \right)^2}.
    \end{equation}
    The final bound follows from the fact that $c\sqrt{\alpha_k} \leq 1$, which follows directly from the hypothesis $0 \leq \alpha_{k+1} \leq \alpha_k - c\alpha_k^{\frac{3}{2}}$.
\end{proof}






% TODO maybe i need the remark before the theorem... to show that the learning rate case and weak case differs from the classical theory!... ie first "we have that (inequalities only using powers of g_k). hence if we have small norm, then... linear convergence. REMARK somewhere about the non-vanishing eta=1=gamma term.


\clearpage

\begin{lemma}
Let the regularization be $\lambda_k = C \sqrt{M\|g_k\|}$ for some $C\geq 1$. If $L$ is $\mu$-strongly convex with $2M$-Lipschitz Hessian and weak gradient edge $\gamma$, and if there exists a $k_0\geq 0$ such that $\|g_{k_0}\| \leq \frac{\mu^2}{4C^2M}\left( \frac{1-\sqrt{1-\gamma^2}}{1+\sqrt{1-\gamma^2}}\right)^2$, then for all $k\geq k_0$ we have that
\begin{align}
    \|g_{k+1}\| \leq \frac{3 + \rho}{4} \|g_k\|
\end{align}
where $\rho = 1 - \eta(1-\sqrt{1-\gamma^2})$. Hence the iterates of \cref{alg:boostingGradRegNewton} converge linearly locally.
\end{lemma}
\begin{proof}
    We have the standard bound
    \begin{align}
        \|g^w_k\| 
        &\leq \|g_k\| + \|g_k^w - g_k\| \\
        &\leq \left( 1 + \sqrt{1-\gamma^2} \right) \|g_k\| \\
        &:= C_\gamma \|g_k\|,
    \end{align}
    and since we assumed $\mu$-strong convexity, we have that
    \begin{align}
        \|f^w_{k+1}\| 
        &= \|(H_k+\lambda I)^{-1} g^w\| \\
        &\leq \frac{1}{\mu}\|g^w\| \\
        &\leq \frac{C_\gamma}{\mu}\|g^w\|.
    \end{align}
    We combine the above with \cref{lemmaUsefulProperties,appendixLemmaGradientGrowth}, specifically \cref{eqGradientRecursionForLater}, to obtain that
    \begin{align}
        \|g_{k+1}\|
        &\leq (1-\eta) \|g_k\| + \eta \|g_k - g_k^w\| +  \eta \lambda_k \|f^w_{k+1}\| + M \eta^2 \|f^w_{k+1}\|^2 \\
        &\leq (1-2\eta + \eta C_\gamma) \|g_k\| + \eta \lambda_k \|f^w_{k+1}\| + M \eta^2 \|f^w_{k+1}\|^2  \\
        &\leq (1-2\eta + \eta C_\gamma) \|g_k\| + \sqrt{C^2 M \eta^2 \frac{C_\gamma^2}{\mu^2}} \|g_k\|^\frac{3}{2} + C^2 M \eta^2 \frac{C_\gamma^2}{\mu^2}\|g_k\|^2.
    \end{align}
    Write $\rho = (1-2\eta + \eta C_\gamma)$ and $K = C^2 M \eta^2 \frac{C_\gamma^2}{\mu^2}$. We seek the condition for a strict decrease $\| g_{k+1} \| < \|g_k\|$. Substituting the above and dividing by $\|g_k\|$ we obtain that
    \begin{align}
        \sqrt{K \|g_k\|} + K\|g_k\| < 1-\rho.
    \end{align}
    Let $\Delta = 1 - \rho = \eta(1 - \sqrt{1-\gamma^2})$. Since $\eta \leq 1$, we observe that $\Delta \leq 1$. Let $y = \sqrt{K\|g_k\|}$. The condition becomes $y + y^2 < \Delta$. A sufficient condition for this to hold is $y \leq \frac{\Delta}{2}$, because this implies $y+y^2 \leq \frac{\Delta}{2} + \frac{\Delta^2}{4} = \frac{\Delta(2+\Delta)}{4}$. Since $\Delta \leq 1$, we have $(2+\Delta) \leq 3$, and thus $y+y^2 \leq \frac{3}{4}\Delta < \Delta$. Expanding the condition $y \leq \frac{\Delta}{2}$ yields
    \begin{align}
        \sqrt{\frac{M \eta^2 C_\gamma^2}{\mu^2} \|g_k\|} \leq \frac{\eta(1 - \sqrt{1-\gamma^2})}{2}.
    \end{align}
    Canceling $\eta$ and rearranging for $\sqrt{\|g_k\|}$, we obtain
    \begin{align}
        \sqrt{\|g_k\|} \leq \frac{\mu}{2\sqrt{M}} \frac{1 - \sqrt{1-\gamma^2}}{C_\gamma} = \frac{\mu}{2\sqrt{M}} \frac{1 - \sqrt{1-\gamma^2}}{1 + \sqrt{1-\gamma^2}}.
    \end{align}
    Squaring both sides gives the bound in the hypothesis. Under this condition, the sum of the higher order terms is bounded by $\frac{3}{4}\Delta \|g_k\| = \frac{3}{4}(1-\rho)\|g_k\|$. Plugging this back into the gradient bound we obtain that
    \begin{align}
        \|g_{k+1}\| \leq \rho \|g_k\| + \frac{3}{4}(1-\rho)\|g_k\| = \frac{3+\rho}{4}\|g_k\|.
    \end{align}
    Since $\rho < 1$, the factor is strictly less than 1, proving linear convergence.
\end{proof}











\clearpage
\section{Divergence of Vanilla Newton on Smooth Convex Losses}\label{app:newton_divergence}
This appendix records explicit smooth convex objectives for which the classical
(unregularized) Newton step may fail to converge globally, even though the loss
satisfies the Lipschitz-Hessian assumption used throughout \cref{secGradientRegularizedNewton} and
\cref{app:GlobalConvergenceProofs}. These examples show that a fixed damping (constant learning rate)
does not, by itself, globalize Newton's method, and motivate the adaptive
gradient regularization used in \cref{alg:boostingGradRegNewton}.



\subsection{Newton, damped Newton, and gradient regularized Newton}

\begin{definition}[Newton-type updates in $\mathbb{R}^d$]
\label{def:newton_updates}
Let $L:\mathbb{R}^d\to\mathbb{R}$ be twice continuously differentiable and assume
$\nabla^2 L(x)\succ 0$ at the current iterate.
\begin{itemize}
  \item The (classical) Newton update is
  \[
    N(x) := x - (\nabla^2 L(x))^{-1}\nabla L(x).
  \]
  \item Given a fixed learning rate $\eta>0$, the damped Newton update is
  \[
    N_\eta(x) := x - \eta(\nabla^2 L(x))^{-1}\nabla L(x).
  \]
\end{itemize}
Assume moreover that $\nabla^2 L$ is $2M$-Lipschitz in operator norm:
\[
  \|\nabla^2 L(u)-\nabla^2 L(v)\|_{\mathrm{op}}\le 2M\|u-v\|
  \qquad \forall u,v\in\mathbb{R}^d.
\]
Then the (exact) gradient regularized Newton (GRN) update is
\[
  \mathrm{GRN}_{M,\eta}(x) := x - \eta\big(\nabla^2L(x)+\lambda(x)I\big)^{-1}\nabla L(x),
  \qquad \lambda(x) := \sqrt{M\|\nabla L(x)\|}.
\]
\end{definition}

\subsection{A canonical 1D counterexample: $L(x)=\sqrt{1+x^2}$}

\begin{proposition}[A smooth convex loss with divergent Newton (1D)]
\label{prop:sqrt_divergence}
Let $L:\mathbb{R}\to\mathbb{R}$ be given by $L(x):=\sqrt{1+x^2}$.
Then $L$ is convex and smooth, and $L''$ is globally Lipschitz.
The classical Newton iteration $x_{k+1}=N(x_k)$ satisfies the exact recursion
\[
  x_{k+1} = -x_k^3.
\]
In particular, $|x_{k+1}|=|x_k|^3$, yielding cubic local convergence for $|x_0|<1$
and divergence for $|x_0|>1$.
\end{proposition}

\begin{proof}
We compute
\[
L'(x)=\frac{x}{\sqrt{1+x^2}},\qquad
L''(x)=\frac{1}{(1+x^2)^{3/2}},\qquad
L^{(3)}(x)=-\frac{3x}{(1+x^2)^{5/2}}.
\]
In particular $L''(x)>0$ for all $x$, hence $L$ is convex. Since $L^{(3)}$ is bounded,
$L''$ is globally Lipschitz. Moreover,
\[
\frac{L'(x)}{L''(x)}
=\frac{x}{\sqrt{1+x^2}}\,(1+x^2)^{3/2}
=x(1+x^2).
\]
Thus the Newton update is
\[
x_{k+1}=x_k-\frac{L'(x_k)}{L''(x_k)}
=x_k-x_k(1+x_k^2)
=-x_k^3.
\]
Finally, $|x_{k+1}|=|x_k|^3$ implies $x_k\to 0$ if $|x_0|<1$ and $|x_k|\to\infty$
if $|x_0|>1$.
\end{proof}

\begin{proposition}[Fixed damping does not globalize Newton]
\label{prop:damped_not_global}
Fix any learning rate $\eta>0$ and consider $x_{k+1}=N_\eta(x_k)$ for
$L(x)=\sqrt{1+x^2}$. Define
\[
  X(\eta) := \sqrt{\max\Big\{0,\frac{3}{\eta}-1\Big\}}.
\]
Then every initialization with $|x_0|\ge X(\eta)$ diverges and satisfies
\[
  |x_{k+1}|\ge 2|x_k| \qquad \forall k\ge 0.
\]
\end{proposition}

\begin{proof}
From the computations above, $L'(x)/L''(x)=x(1+x^2)$, hence
\[
x_{k+1}=x_k-\eta x_k(1+x_k^2)=x_k\big(1-\eta(1+x_k^2)\big).
\]
If $|x_k|\ge X(\eta)$, then by definition $\eta(1+x_k^2)\ge 3$, so
$|1-\eta(1+x_k^2)|\ge 2$ and therefore $|x_{k+1}|\ge 2|x_k|$.
Induction yields divergence.
\end{proof}

\subsection{A drift-based generator for smooth convex Newton counterexamples}
\label{app:drift_generator}

The previous example is not accidental. We now record a general construction
principle: prescribing a one-dimensional \emph{Newton drift} $d(x)$ determines a
convex objective $L$ whose Newton iteration is exactly $x_{k+1}=x_k-\tilde d(x_k)$,
where $\tilde d$ is the odd extension of $d$. This yields an infinite family of
smooth convex losses on which classical Newton diverges.

\begin{theorem}[Lipschitz-Hessian drift-to-loss generator for 1D Newton dynamics]
\label{thm:drift_to_loss_lipschitz}
Let $d:[0,\infty)\to[0,\infty)$ be continuous and satisfy:
{\renewcommand{\labelenumi}{(D\arabic{enumi})}%
\begin{enumerate}
\setcounter{enumi}{-1}
\item $d(0)=0$ and $d(x)>0$ for all $x>0$;
\item there exist $\delta>0$ and $C_0<\infty$ such that
\[
|d(x)-x|\le C_0x^2 \qquad \forall x\in[0,\delta].
\]
\end{enumerate}}

Define, for $x>0$,
\[
\Phi(x):=\int_1^x \frac{ds}{d(s)},\qquad
L_+'(x):=a\,e^{\Phi(x)}\ (a>0),\qquad
L_+(x):=\int_0^x L_+'(t)\,dt,
\]
and extend evenly by $L(x):=L_+(|x|)$. Let $\tilde d:\mathbb{R}\to\mathbb{R}$ be the odd extension
$\tilde d(x):=\mathrm{sign}(x)\,d(|x|)$ with $\tilde d(0)=0$. Then:
{\renewcommand{\labelenumi}{(\roman{enumi})}%
\begin{enumerate}
\item $L$ is even, convex, and $C^2(\mathbb{R})$, with $L''(x)>0$ for $x\neq 0$.
Moreover, for $x>0$,
\[
L'(x)=L_+'(x),\qquad
L''(x)=\frac{L_+'(x)}{d(x)}.
\]
In particular, $L'$ is odd and $L''$ is even.

\item For every $x\neq 0$,
\[
\frac{L'(x)}{L''(x)}=\tilde d(x),
\]
and therefore the classical Newton update for $L$ is exactly
\[
N(x)=x-\frac{L'(x)}{L''(x)}=x-\tilde d(x).
\]

\item If there exist $X_0>0$ and $\beta>1$ such that
\[
d(x)\ge (1+\beta)x \qquad \forall x\ge X_0,
\]
then every Newton iterate initialized with $|x_0|\ge X_0$ diverges and satisfies
\[
|x_{k+1}|\ge \beta|x_k| \qquad \forall k\ge 0.
\]

\item \textbf{(Globally Lipschitz Hessian.)}
Assume in addition that $d\in C^1((0,\infty))$ and that
\begin{equation}
\label{eq:D2_ratio}
C_1 \;:=\;\sup_{x>0}\frac{|d'(x)-1|}{d(x)}\;<\;\infty.
\end{equation}
Assume moreover that there exist $\kappa>1$ and $X_1>0$ such that
\begin{equation}
\label{eq:D3_linear_lower}
d(x)\ge \kappa x \qquad \forall x\ge X_1.
\end{equation}
(For instance, \cref{eq:D3_linear_lower} holds whenever the hypothesis in (iii) holds.)
Then $L''$ is globally Lipschitz on $\mathbb{R}$, i.e.\ there exists $M<\infty$ such that
\[
|L''(u)-L''(v)|\le 2M|u-v|\qquad \forall u,v\in\mathbb{R}.
\]
More precisely, for $x>0$ we have
\begin{equation}
\label{eq:L3_formula}
L^{(3)}(x)=\frac{L_+'(x)}{d(x)^2}\,\bigl(1-d'(x)\bigr)
= L''(x)\,\frac{1-d'(x)}{d(x)},
\end{equation}
and hence $\sup_{x>0}|L^{(3)}(x)|<\infty$, so one may take
$2M=\sup_{x\neq 0}|L^{(3)}(x)|$.
\end{enumerate}}
\end{theorem}

\begin{proof}
We split the argument into four parts.

\paragraph{Part A: construction and smoothness on $(0,\infty)$.}
By (D0), the map $s\mapsto 1/d(s)$ is continuous on compact subsets of $(0,\infty)$.
Hence $\Phi\in C^1((0,\infty))$ with $\Phi'(x)=1/d(x)$.
Therefore $L_+'(x)=a e^{\Phi(x)}$ is $C^1((0,\infty))$ and
\[
(L_+')'(x)=a e^{\Phi(x)}\Phi'(x)=\frac{L_+'(x)}{d(x)}.
\]
Thus $L_+\in C^2((0,\infty))$, with
\[
L_+''(x)=(L_+')'(x)=\frac{L_+'(x)}{d(x)}>0 \qquad (x>0),
\]
so $L_+$ is strictly convex on $(0,\infty)$.

\paragraph{Part B: behavior at $0$ and $C^2$ even extension.}
Write $d(x)=x(1+\varepsilon(x))$ on $(0,\delta]$, where by (D1) we have
$|\varepsilon(x)|\le C_0x$, hence $\varepsilon(x)\to 0$ as $x\downarrow 0$.
Then on $(0,\delta]$,
\[
\frac{1}{d(x)}=\frac{1}{x}\cdot \frac{1}{1+\varepsilon(x)}
=\frac{1}{x}+g(x),
\]
where $g$ is bounded on $(0,\delta]$ (since $(1+\varepsilon)^{-1}=1+O(\varepsilon)$).
Integrating gives
\[
\Phi(x)=\int_1^x \frac{ds}{d(s)}
=\int_1^x \frac{ds}{s}+\int_1^x g(s)\,ds
=\log x + C_\ast + o(1)\qquad (x\downarrow 0)
\]
for some finite constant $C_\ast$.
Exponentiating yields
\[
L_+'(x)=a e^{\Phi(x)}=a e^{C_\ast}\,x(1+o(1))\qquad (x\downarrow 0),
\]
so $L_+'(x)\to 0$ and $L_+'(x)/x\to a e^{C_\ast}\in(0,\infty)$.
Since $d(x)=x(1+o(1))$, we obtain
\[
L_+''(x)=\frac{L_+'(x)}{d(x)}\to a e^{C_\ast}\in(0,\infty)\qquad (x\downarrow 0).
\]
Define $L_+(0):=0$, $L_+'(0):=0$, and $L_+''(0):=\lim_{x\downarrow 0}L_+''(x)$; then
$L_+\in C^2([0,\infty))$.

Now define $L(x)=L_+(|x|)$. This $L$ is even and $C^2(\mathbb{R})$ with
$L'(x)$ odd and $L''(x)$ even, and for $x>0$ we have
$L'(x)=L_+'(x)$ and $L''(x)=L_+''(x)=L_+'(x)/d(x)$.
Since $L''(x)\ge 0$ everywhere and $L''(x)>0$ for $x\neq 0$, $L$ is convex.
This proves (i).

\paragraph{Part C: drift identity and divergence.}
For $x>0$, by (i) we have $L'(x)=L_+'(x)$ and $L''(x)=L_+'(x)/d(x)$, hence
\[
\frac{L'(x)}{L''(x)}=d(x).
\]
By evenness of $L''$ and oddness of $L'$, the ratio $L'/L''$ is odd, so for $x<0$,
\[
\frac{L'(x)}{L''(x)}=-\frac{L'(|x|)}{L''(|x|)}=-d(|x|)=\tilde d(x),
\]
proving (ii). The divergence statement (iii) follows exactly as in \cref{thm:drift_to_loss_lipschitz}:
if $x\ge X_0$, then $N(x)=x-d(x)\le -\beta x$, hence $|N(x)|\ge \beta|x|$, and similarly for $x\le -X_0$.

\paragraph{Part D: globally Lipschitz Hessian.}
Assume \cref{eq:D2_ratio,eq:D3_linear_lower}.
For $x>0$, differentiating $L''(x)=L_+'(x)/d(x)$ and using $(L_+')'(x)=L_+'(x)/d(x)$
from Part~A gives \cref{eq:L3_formula}. Hence, for $x>0$,
\[
|L^{(3)}(x)| \le \frac{L_+'(x)}{d(x)^2}\,|1-d'(x)|
\le C_1 \frac{L_+'(x)}{d(x)} = C_1 L''(x).
\]
It remains to show $\sup_{x>0}L''(x)<\infty$.
By Part~B, $L''$ extends continuously to $x=0$ with $0<L''(0)<\infty$.
On $[0,X_1]$, continuity gives $\sup_{[0,X_1]}L''<\infty$.
For $x\ge X_1$, using \cref{eq:D3_linear_lower},
\[
\Phi(x)=\Phi(X_1)+\int_{X_1}^x \frac{ds}{d(s)}
\le \Phi(X_1)+\int_{X_1}^x \frac{ds}{\kappa s}
=\Phi(X_1)+\frac{1}{\kappa}\log\frac{x}{X_1},
\]
so $L_+'(x)=a e^{\Phi(x)}\le C\,x^{1/\kappa}$ for some $C$.
Therefore,
\[
L''(x)=\frac{L_+'(x)}{d(x)}\le \frac{C x^{1/\kappa}}{\kappa x}
= \frac{C}{\kappa}x^{\frac{1}{\kappa}-1},
\]
which is bounded for $x\ge X_1$ since $\kappa>1$.
Thus $\sup_{x>0}|L^{(3)}(x)|<\infty$, and since $L''$ is even, $L''$ is globally Lipschitz on $\mathbb{R}$.
\end{proof}

\subsection{Explicit Counterexamples }

\begin{corollary}[Elementary closed-form losses from explicit drifts]
\label{cor:elementary_losses_from_drifts}
Applying \cref{thm:drift_to_loss_lipschitz} with the following choices of $d$
(and the corresponding choice of $a$) yields the stated closed-form losses:
\begin{enumerate}
\item If $d(x)=x(1+x)$ and $a=C/2$, then
\[
L(x)=C\bigl(|x|-\ln(1+|x|)\bigr).
\]

\item If $d(x)=x(1+x^2)$ and $a=C/\sqrt{2}$, then
\[
L(x)=C\bigl(\sqrt{1+x^2}-1\bigr).
\]

\item Fix an integer $m\ge 3$. If
\[
d(x)=\frac{1+x}{m-1}\Bigl((1+x)^{m-1}-1\Bigr)
\qquad\text{and}\qquad
a=C\Bigl(1-2^{-(m-1)}\Bigr),
\]
then
\[
L(x)=C|x|-\frac{C}{m-2}\Bigl(1-(1+|x|)^{-(m-2)}\Bigr).
\]

\item If $d(x)=(1+x^2)\arctan x$ and $a=C(\pi/4)$, then
\[
L(x)=C\Bigl(|x|\arctan|x|-\tfrac12\ln(1+x^2)\Bigr).
\]
\end{enumerate}
In each case, $L$ is even, convex, $C^2$ with globally Lipschitz Hessian, and the
classical Newton method diverges from sufficiently large initializations.
\end{corollary}

\begin{proof}
For each item, one checks that the stated $L$ has $L'(x)/L''(x)=\tilde d(x)$ for $x\neq 0$,
and that $L_+'(1)=a$ matches the choice of $a$ above; therefore \cref{thm:drift_to_loss_lipschitz}
constructs exactly the stated $L$. The regularity and divergence conclusions follow from the theorem. We indicate this in the case of $d(x)=x+x^3=x(1+x^2)$. We compute, for $x>0$,
\[
\Phi(x)=\int_1^x\frac{ds}{s(1+s^2)}
=\Big[\log s-\tfrac12\log(1+s^2)\Big]_{s=1}^{s=x}
=\log x-\tfrac12\log(1+x^2)+\tfrac12\log 2.
\]
Therefore
\[
e^{\Phi(x)}=\sqrt{2}\,\frac{x}{\sqrt{1+x^2}},
\qquad
L_+'(x)=a e^{\Phi(x)}=\frac{x}{\sqrt{1+x^2}}
\quad\text{when } a=2^{-1/2}.
\]
Integrating from $0$ to $x$ and using $L(0)=0$ yields $L(x)=\sqrt{1+x^2}-1$.
Finally, by \cref{thm:drift_to_loss_lipschitz}(ii),
\[
x_{k+1}=x_k-\tilde d(x_k)=x_k-(x_k+x_k^3)=-x_k^3,
\]
as claimed.
\end{proof}


\subsection{A boosting/Hilbert-space realization}

\begin{proposition}[A separable empirical-risk objective with divergent Newton updates]
\label{prop:boosting_realization}
Let $\mathcal{H}=L^2(\hat\nu_N)\cong \mathbb{R}^N$ be the space of predictions on the training points,
and consider the empirical risk
\[
  L(F) := \frac1N\sum_{i=1}^N \sqrt{1 + F(x_i)^2}.
\]
Then $L$ is convex and has $2M$-Lipschitz Hessian for some $M<\infty$, and the exact Newton update satisfies, coordinate-wise,
\[
  F_{k+1}(x_i) = -F_k(x_i)^3,\qquad i=1,\dots,N.
\]
In particular, if $|F_0(x_i)|>1$ for some $i$, then $\|F_k\|\to\infty$.
\end{proposition}

\begin{proof}
The objective is separable across coordinates in $\mathbb{R}^N$. The coordinate-wise gradient and
Hessian are those of $L(x)=\sqrt{1+x^2}$, so \cref{prop:sqrt_divergence} applies componentwise.
Lipschitzness of the Hessian follows because the scalar third derivative is bounded, hence the diagonal
Hessian map is Lipschitz in operator norm.
\end{proof}

\subsection{Why gradient regularization fixes it}

\begin{proposition}[GRN converges on the same examples]
\label{prop:grn_converges}
Consider any convex objective $L:\mathbb{R}^d\to\mathbb{R}$ with $2M$-Lipschitz Hessian and finite sublevel sets. Then the exact gradient regularized Newton method in \cref{def:newton_updates}
converges globally with function-value rate $L(x_k)-\min_x L(x)=O(1/k^2)$ for any fixed $\eta\in(0,1]$.
In particular, this applies to $L(x)=\sqrt{1+x^2}$ and to all losses generated by
\cref{thm:drift_to_loss_lipschitz}.
\end{proposition}

\begin{proof}
This is a direct specialization of \cref{thm:global_rate_II} / \cref{app:GlobalConvergenceProofs},
since the weak-step becomes exact and the weak-gradient edge is $1$.
\end{proof}

\newpage

\section{Newton Boosting for Hessian-Dominated Losses}\label{app:HessianDominatedLoss}



This appendix proves the claims stated around \cref{defHessianDominatedLoss}: binary and
categorical cross entropy are Hessian-dominated with constant $c=1$, and
pointwise Hessian-dominance implies Hessian-dominance of the empirical risk.

\subsection{A useful log inequality}

\begin{lemma}[A log inequality]
\label{lem:log_ineq}
For all $p\in(0,1]$,
\begin{equation}
\label{eq:log_ineq_1}
\frac{1-p}{p} \;\ge\; -\log p.
\end{equation}
Equivalently, for all $p\in[0,1)$,
\begin{equation}
\label{eq:log_ineq_2}
\frac{p}{1-p} \;\ge\; -\log(1-p).
\end{equation}
\end{lemma}

\begin{proof}
Define $\phi(p)=\frac{1}{p}-1+\log p$ on $(0,1]$. Then
\[
\phi'(p)=-\frac{1}{p^2}+\frac{1}{p}=\frac{p-1}{p^2}\le 0,
\]
so $\phi$ is non-increasing. Since $\phi(1)=0$, we have $\phi(p)\ge 0$ for all
$p\in(0,1]$, which is \eqref{eq:log_ineq_1}. Substituting $q=1-p$ gives
\eqref{eq:log_ineq_2}.
\end{proof}


\subsection{Binary cross entropy is Hessian-dominated}

\begin{proposition}[Binary cross entropy is Hessian-dominated with $c=1$]
\label{prop:binary_ce_hessdom}
Let $y\in\{0,1\}$ and $u\in\mathbb{R}$, and set $p=\sigma(u)=(1+e^{-u})^{-1}$.
Consider the binary cross entropy
\[
l(u,y) \;=\; -\Big(y\log p + (1-y)\log(1-p)\Big).
\]
Let $g(u,y)=\partial_u l(u,y)$ and $h(u)=\partial_{uu}l(u,y)$.
Then the (scalar) Newton step $f(u,y)=-h(u)^{-1}g(u,y)$ satisfies
\begin{equation}
\label{eq:binary_pointwise_dom}
\|f(u,y)\|_{h(u)}^2 := f(u,y)^2\,h(u)
= \frac{g(u,y)^2}{h(u)}
\;\ge\; l(u,y).
\end{equation}
Hence binary cross entropy is pointwise Hessian-dominated with constant $c=1$.
\end{proposition}

\begin{proof}
A direct computation gives
\[
g(u,y)=p-y,\qquad h(u)=p(1-p)>0.
\]
Thus $\|f\|_{h}^2 = g^2/h$. If $y=1$, then $g=p-1=-(1-p)$, so
\[
\frac{g^2}{h}=\frac{(1-p)^2}{p(1-p)}=\frac{1-p}{p}.
\]
Also $l(u,1)=-\log p$, and \cref{lem:log_ineq} gives
$\frac{1-p}{p}\ge -\log p=l(u,1)$. If $y=0$, then $g=p$, so
\[
\frac{g^2}{h}=\frac{p^2}{p(1-p)}=\frac{p}{1-p}.
\]
Also $l(u,0)=-\log(1-p)$, and \cref{lem:log_ineq} gives
$\frac{p}{1-p}\ge -\log(1-p)=l(u,0)$. This proves \eqref{eq:binary_pointwise_dom}.
\end{proof}

\begin{lemma}[Softmax Hessian is positive definite on the zero-sum subspace]
\label{lem:softmax_hess_posdef}
Let $p\in\mathbb{R}^K$ satisfy $p_k>0$ and $\sum_{k=1}^K p_k=1$, and define
\[
H \;:=\; \mathrm{diag}(p)-pp^\top.
\]
Let $V:=\{v\in\mathbb{R}^K:\mathbf{1}^\top v=0\}$. Then $H$ maps $V$ to $V$ and is
positive definite on $V$. In particular, the restriction $H:V\to V$ is invertible.
\end{lemma}

\begin{proof}
First, for any $v\in\mathbb{R}^K$,
\[
\mathbf{1}^\top H v
= \mathbf{1}^\top \mathrm{diag}(p)v - \mathbf{1}^\top p (p^\top v)
= p^\top v - (p^\top v) = 0,
\]
so $Hv\in V$ whenever $v\in V$. Next, for any $v\in\mathbb{R}^K$,
\[
v^\top H v
= \sum_{k=1}^K p_k v_k^2 - \Big(\sum_{k=1}^K p_k v_k\Big)^2
= \sum_{k=1}^K p_k\big(v_k-\sum_{j=1}^K p_j v_j\big)^2 \;\ge\; 0.
\]
Equality holds iff $v_k$ is constant in $k$, i.e.\ $v\in\mathrm{span}\{\mathbf{1}\}$.
Intersecting with $V$ leaves only $v=0$, hence $v^\top H v>0$ for all nonzero
$v\in V$. Therefore $H$ is positive definite on $V$ and invertible as $H:V\to V$.
\end{proof}

\begin{proposition}[Categorical cross entropy is Hessian-dominated with $c=1$]
\label{prop:categorical_ce_hessdom}
Let $u\in\mathbb{R}^K$, $p=\mathrm{softmax}(u)$, and let $y=e_t$ be a one-hot label.
Consider categorical cross entropy
\[
l(u,y) \;=\; -\sum_{k=1}^K y_k\log p_k \;=\; -\log p_t.
\]
Let $g(u,y)=\nabla_u l(u,y)=p-y$ and $H(u)=\nabla^2_{uu}l(u,y)=\mathrm{diag}(p)-pp^\top$.
Work on the subspace $V=\{v:\mathbf{1}^\top v=0\}$; then $g(u,y)\in V$ and by
\cref{lem:softmax_hess_posdef} the Newton direction $f\in V$ solving
\begin{equation}
\label{eq:cat_newton}
H(u)f \;=\; y-p
\end{equation}
is unique. Moreover,
\begin{equation}
\label{eq:cat_decrement}
\|f\|_{H(u)}^2 := f^\top H(u)f
\;=\; \frac{1-p_t}{p_t}
\;\ge\; -\log p_t \;=\; l(u,y).
\end{equation}
Hence categorical cross entropy is pointwise Hessian-dominated with constant $c=1$.
\end{proposition}

\begin{proof}
Let $H=\mathrm{diag}(p)-pp^\top$. Define $s:=p^\top f$. Expanding \eqref{eq:cat_newton} yields
\[
\mathrm{diag}(p)f - p(p^\top f)=y-p
\quad\Longleftrightarrow\quad
p_k(f_k-s)=y_k-p_k\qquad \forall k,
\]
so (since $p_k>0$) we have
\[
f_k = s + \frac{y_k}{p_k} - 1.
\]
With $y=e_t$, this becomes
\[
f_t = s + \frac{1}{p_t}-1,\qquad f_k=s-1\quad(k\neq t).
\]
Using symmetry of $H$ and \eqref{eq:cat_newton},
\[
\|f\|_H^2=f^\top H f = f^\top (y-p).
\]
Compute
\begin{align*}
f^\top (y-p)
&= f_t(1-p_t) + \sum_{k\neq t} f_k(0-p_k) \\
&= f_t(1-p_t) - (s-1)\sum_{k\neq t}p_k \\
&= \big(f_t-(s-1)\big)(1-p_t)
= \left(\left(s+\frac{1}{p_t}-1\right)-(s-1)\right)(1-p_t) \\
&= \frac{1-p_t}{p_t}.
\end{align*}
Finally, \cref{lem:log_ineq} gives $\frac{1-p_t}{p_t}\ge -\log p_t=l(u,y)$,
proving \eqref{eq:cat_decrement}.
\end{proof}

\begin{remark}[Relation to Assumption~3.1 (uniqueness of Newton step)]
\label{rem:cat_uniqueness}
On $\mathbb{R}^K$ one has $H(u)\mathbf{1}=0$ (softmax is shift-invariant), so the
Newton equation is not uniquely solvable without fixing a gauge. Restricting to the
zero-sum subspace $V=\{v:\mathbf{1}^\top v=0\}$ (or equivalently fixing one logit
coordinate) yields uniqueness and matches Assumption~3.1.
The decrement $\|f\|_{H(u)}^2$ is invariant to the choice of representative.
\end{remark}

\begin{remark}
Binary cross entropy corresponds to the special case $K=2$ of categorical
cross entropy (after identifying the scalar logit with the difference of the
two softmax logits). Thus \cref{prop:categorical_ce_hessdom}
strictly generalizes \cref{prop:binary_ce_hessdom}; the latter is
included separately for clarity and because it avoids the gauge-fixing
required in the multiclass setting.
\end{remark}

\subsection{From pointwise to empirical Hessian-dominance}

\begin{lemma}[Pointwise $\Rightarrow$ empirical Hessian-dominance]
\label{lem:pointwise_to_empirical_hessdom}
Let $\{(x_i,y_i)\}_{i=1}^N$ be a dataset and define the empirical risk
\[
L(F)=\frac{1}{N}\sum_{i=1}^N l(F(x_i),y_i),
\]
with $F(x_i)\in\mathbb{R}^K$ (take $K=1$ for binary). Work in the empirical Hilbert
space $H=L^2(\hat\nu_N)$ with inner product
\[
\langle A,B\rangle = \frac{1}{N}\sum_{i=1}^N \langle A(x_i),B(x_i)\rangle_{\mathbb{R}^K}.
\]
Fix $F\in H$ and write $u_i:=F(x_i)$. Let
\[
g_{(i)} := \partial_1 l(u_i,y_i)\in\mathbb{R}^K,\qquad
h_{(i)} := \partial_{11} l(u_i,y_i)\in\mathbb{R}^{K\times K}.
\]
Assume there exists $c>0$ such that for all $i$,
\begin{equation}
\label{eq:pointwise_dom_assump}
g_{(i)}^\top h_{(i)}^{-1} g_{(i)} \;\ge\; c\,l(u_i,y_i),
\end{equation}
where $h_{(i)}^{-1}$ is taken on the relevant space (for categorical CE one may restrict
to $V$ pointwise as in \cref{prop:categorical_ce_hessdom}).
Let $g=\nabla L(F)$ and $H=\nabla^2 L(F)$, and let $f=-H^{-1}g$ be the exact Newton
direction (Assumption~3.1). Then
\[
\|f\|_{H}^2 \;\ge\; c\,L(F).
\]
In particular, $L$ is Hessian-dominated in the sense of \cref{defHessianDominatedLoss} with constant $c$.
\end{lemma}

\begin{proof}
In the empirical setting, the gradient and Hessian act pointwise:
\[
g(x_i)=g_{(i)},\qquad (H[f])(x_i)=h_{(i)} f(x_i).
\]
Thus the Newton equation $H[f]=-g$ decouples over samples:
\[
h_{(i)} f(x_i)=-g_{(i)} \quad\Rightarrow\quad f(x_i)=-h_{(i)}^{-1}g_{(i)}.
\]
Therefore, using \cref{defCosineAngle},
\begin{align*}
\|f\|_{H}^2
&= \langle f,H[f]\rangle
= \frac{1}{N}\sum_{i=1}^N f(x_i)^\top h_{(i)} f(x_i)
= \frac{1}{N}\sum_{i=1}^N g_{(i)}^\top h_{(i)}^{-1} g_{(i)}.
\end{align*}
Applying \eqref{eq:pointwise_dom_assump} term-by-term yields
\[
\|f\|_{H}^2
\ge \frac{c}{N}\sum_{i=1}^N l(u_i,y_i)
= c\,L(F),
\]
which is exactly \cref{defHessianDominatedLoss}.
\end{proof}

\begin{corollary}[Empirical cross entropy risks are Hessian-dominated with $c=1$]
\label{cor:empirical_ce_hessdom}
Let $L(F)=\frac{1}{N}\sum_{i=1}^N l(F(x_i),y_i)$ where $l$ is either binary cross entropy
(\cref{prop:binary_ce_hessdom}) or categorical cross entropy
(\cref{prop:categorical_ce_hessdom}). Then $L$ is Hessian-dominated with $c=1$
(in the sense of \cref{defHessianDominatedLoss}, with the gauge-fix of \cref{rem:cat_uniqueness}
in the categorical case).
\end{corollary}














\end{document}


% This document was modified from the file originally made available by
% Pat Langley and Andrea Danyluk for ICML-2K. This version was created
% by Iain Murray in 2018, and modified by Alexandre Bouchard in
% 2019 and 2021 and by Csaba Szepesvari, Gang Niu and Sivan Sabato in 2022.
% Modified again in 2023 and 2024 by Sivan Sabato and Jonathan Scarlett.
% Previous contributors include Dan Roy, Lise Getoor and Tobias
% Scheffer, which was slightly modified from the 2010 version by
% Thorsten Joachims & Johannes Fuernkranz, slightly modified from the
% 2009 version by Kiri Wagstaff and Sam Roweis's 2008 version, which is
% slightly modified from Prasad Tadepalli's 2007 version which is a
% lightly changed version of the previous year's version by Andrew
% Moore, which was in turn edited from those of Kristian Kersting and
% Codrina Lauth. Alex Smola contributed to the algorithmic style files.









































